% papers/wording/expected.tex                                      -*-LaTeX-*-
% SPDX-License-Identifier: Apache-2.0 WITH LLVM-exception
%
% Generated from the annotated headers in include/beman/expected/ via specgen.
% Do not edit by hand: change a header's //! docblocks and run `make wording`
% instead. This comment comes from papers/wording/preamble.tex, which is the
% one part of this file a person writes.
%
% --base-section-depth 2 puts each \rSec marker at the level the draft
% writes it at: [expected.unexpected] and its siblings are \rSec2, their
% subclauses \rSec3, numbering 22.8.3, 22.8.3.1, ... just as in
% source/utilities.tex. This file is therefore everything that sits
% *inside* the existing \rSec1[expected]{Expected objects} in the draft's
% source/utilities.tex, in clause order, ready to replace the current
% [expected.unexpected] through [expected.void] subclauses and add the new
% [expected.ref] one after them -- no \rSec1[expected] wrapper and no
% \input directives. [expected.general] and [expected.syn] are prose, not
% generated from any one declaration; see papers/expected-new.tex.

\rSec2[expected.unexpected]{Class template unexpected}

\rSec3[expected.un.general]{General}

\begin{codeblock}
namespace std {
  template<class E>
  class @\libglobal{unexpected}@ {
  public:
    constexpr unexpected(const unexpected&) = default;
    constexpr unexpected(unexpected&&) = default;

    template<class Err = E>
    constexpr explicit unexpected(Err&& e) noexcept(is_nothrow_constructible_v<E, Err>);

    template<class... Args>
    constexpr explicit unexpected(in_place_t, Args&&... args) noexcept(
        is_nothrow_constructible_v<E, Args...>);

    template<class U, class... Args>
    constexpr explicit unexpected(
        in_place_t, initializer_list<U> il,
        Args&&... args) noexcept(is_nothrow_constructible_v<E, initializer_list<U>&,
                                                            Args...>);

    constexpr unexpected& operator=(const unexpected&) = default;
    constexpr unexpected& operator=(unexpected&&) = default;

    constexpr const E& error() const& noexcept;
    constexpr E& error() & noexcept;
    constexpr const E&& error() const&& noexcept;
    constexpr E&& error() && noexcept;

    constexpr void swap(unexpected& other) noexcept(is_nothrow_swappable_v<E>);

    template<class E2>
    friend constexpr bool operator==(const unexpected& x, const unexpected<E2>& y);

    friend constexpr void swap(unexpected& x,
                               unexpected& y) noexcept(noexcept(x.swap(y)));

  private:
    E @\exposid{unex}@; // \expos
  };
}
\end{codeblock}

\pnum
\mandates
A program that instantiates the definition of \tcode{unexpected} for a non-object type other than an lvalue reference type, an array type, a specialization of \tcode{unexpected}, or a cv-qualified type is ill-formed.

\pnum
\remarks
Subclause \iref{expected.unexpected} describes the class template \tcode{unexpected} that represents unexpected objects stored in \tcode{expected} objects.

\rSec3[expected.un.cons]{Constructors}

\indexlibraryctor{unexpected}%
\begin{itemdecl}
template<class Err = E>
  constexpr explicit unexpected(Err&& e) noexcept(is_nothrow_constructible_v<E, Err>);
\end{itemdecl}

\begin{itemdescr}
\pnum
\constraints
\begin{itemize}
\item \tcode{is_same_v<remove_cvref_t<Err>, unexpected>} is \tcode{false}; and
\item \tcode{is_same_v<remove_cvref_t<Err>, in_place_t>} is \tcode{false}; and
\item \tcode{is_constructible_v<E, Err>} is \tcode{true}.
\end{itemize}

\pnum
\effects
Direct-non-list-initializes \exposid{unex} with \tcode{std::forward<Err>(e)}.

\pnum
\throws
Any exception thrown by the initialization of \exposid{unex}.
\end{itemdescr}

\indexlibraryctor{unexpected}%
\begin{itemdecl}
template<class... Args>
  constexpr explicit unexpected(in_place_t, Args&&... args) noexcept(is_nothrow_constructible_v<E, Args...>);
\end{itemdecl}

\begin{itemdescr}
\pnum
\constraints
\tcode{is_constructible_v<E, Args...>} is \tcode{true}.

\pnum
\effects
Direct-non-list-initializes \exposid{unex} with \tcode{std::forward<Args>(args)...}.

\pnum
\throws
Any exception thrown by the initialization of \exposid{unex}.
\end{itemdescr}

\indexlibraryctor{unexpected}%
\begin{itemdecl}
template<class U, class... Args>
  constexpr explicit unexpected(in_place_t, initializer_list<U> il,
                                Args&&... args) noexcept(is_nothrow_constructible_v<E, initializer_list<U>&, Args...>);
\end{itemdecl}

\begin{itemdescr}
\pnum
\constraints
\tcode{is_constructible_v<E, initializer_list<U>\&, Args...>} is \tcode{true}.

\pnum
\effects
Direct-non-list-initializes \exposid{unex} with \tcode{il, std::forward<Args>(args)...}.

\pnum
\throws
Any exception thrown by the initialization of \exposid{unex}.
\end{itemdescr}

\rSec3[expected.un.obs]{Observers}

\indexlibrarymember{error}{unexpected}%
\begin{itemdecl}
constexpr const E& error() const& noexcept;
constexpr E& error() & noexcept;
\end{itemdecl}

\begin{itemdescr}
\pnum
\returns
\exposid{unex}.
\end{itemdescr}

\indexlibrarymember{error}{unexpected}%
\begin{itemdecl}
constexpr const E&& error() const&& noexcept;
constexpr E&& error() && noexcept;
\end{itemdecl}

\begin{itemdescr}
\pnum
\returns
\tcode{std::move(\exposid{unex})}.
\end{itemdescr}

\rSec3[expected.un.swap]{Swap}

\indexlibrarymember{swap}{unexpected}%
\begin{itemdecl}
constexpr void swap(unexpected& other) noexcept(is_nothrow_swappable_v<E>);
\end{itemdecl}

\begin{itemdescr}
\pnum
\mandates
\tcode{is_swappable_v<E>} is \tcode{true}.

\pnum
\effects
Equivalent to: \tcode{using std::swap; swap(\exposid{unex}, other.\exposid{unex});}
\end{itemdescr}

\indexlibraryglobal{swap}%
\begin{itemdecl}
friend constexpr void swap(unexpected& x, unexpected& y) noexcept(noexcept(x.swap(y)));
\end{itemdecl}

\begin{itemdescr}
\pnum
\constraints
\tcode{is_swappable_v<E>} is \tcode{true}.

\pnum
\effects
Equivalent to \tcode{x.swap(y)}.
\end{itemdescr}

\rSec3[expected.un.eq]{Equality operator}

\indexlibraryglobal{operator==}%
\begin{itemdecl}
template<class E2>
  friend constexpr bool operator==(const unexpected& x, const unexpected<E2>& y);
\end{itemdecl}

\begin{itemdescr}
\pnum
\mandates
The expression \tcode{x.error() == y.error()} is well-formed and its result is convertible to \tcode{bool}.

\pnum
\returns
\tcode{x.error() == y.error()}.
\end{itemdescr}

\rSec3[expected.un.ref]{Partial specialization unexpected<E\&>}

\begin{codeblock}
namespace std {
  template<class E>
  class unexpected<E&> {
  public:
    constexpr unexpected(const unexpected&) = default;
    constexpr unexpected(unexpected&&) = default;
    template<class G = E>
      constexpr explicit unexpected(G&&) noexcept;
    template<class G = E>
      constexpr explicit unexpected(in_place_t, G&&) noexcept;

    constexpr unexpected& operator=(const unexpected&) = default;
    constexpr unexpected& operator=(unexpected&&) = default;

    constexpr E& error() const noexcept;

    constexpr void swap(unexpected& other) noexcept;

    template<class E2>
      friend constexpr bool operator==(const unexpected&, const unexpected<E2>&);

    friend constexpr void swap(unexpected& x, unexpected& y) noexcept;

  private:
    E* unex;              // \expos
  };
}
\end{codeblock}

\pnum
\mandates
A program that instantiates the definition of \tcode{unexpected<E\&>} for a referenced type that is not an object type, an array type, or a specialization of \tcode{unexpected} is ill-formed.

\pnum
\remarks
An object of type \tcode{unexpected<E\&>} holds a pointer to an object of type \tcode{E}. The referenced object is not owned by the \tcode{unexpected<E\&>} object. Unlike the primary template, \tcode{E} may be a cv-qualified type.

\indexlibraryctor{unexpected}%
\begin{itemdecl}
template<class G = E>
  constexpr explicit unexpected(G&& e) noexcept;
\end{itemdecl}

\begin{itemdescr}
\pnum
\constraints
\begin{itemize}
\item \tcode{is_same_v<remove_cvref_t<G>, unexpected>} is \tcode{false},
\item \tcode{is_same_v<remove_cvref_t<G>, in_place_t>} is \tcode{false},
\item \tcode{is_constructible_v<E\&, G>} is \tcode{true}, and
\item \tcode{reference_constructs_from_temporary_v<E\&, G>} is \tcode{false}.
\end{itemize}

\pnum
\effects
Initializes \exposid{unex} with \tcode{addressof(static_cast<E\&>(std::forward<G>(e)))}.

\pnum
\remarks
A constructor for which \tcode{reference_constructs_from_temporary_v<E\&, G>} is \tcode{true} — one that would bind \tcode{E\&} to a temporary — is defined as deleted.
\end{itemdescr}

\indexlibraryctor{unexpected}%
\begin{itemdecl}
template<class G = E>
  constexpr explicit unexpected(in_place_t, G&& e) noexcept;
\end{itemdecl}

\begin{itemdescr}
\pnum
\constraints
\tcode{is_constructible_v<E\&, G>} is \tcode{true} and \tcode{reference_constructs_from_temporary_v<E\&, G>} is \tcode{false}.

\pnum
\effects
Initializes \exposid{unex} with \tcode{addressof(static_cast<E\&>(std::forward<G>(e)))}.
\end{itemdescr}

\indexlibrarymember{error}{unexpected}%
\begin{itemdecl}
constexpr E& error() const noexcept;
\end{itemdecl}

\begin{itemdescr}
\pnum
\returns
\tcode{*\exposid{unex}}.

\pnum
\remarks
The reference returned is not \tcode{const}-qualified even when \tcode{*this} is \tcode{const}; the constness of the \tcode{unexpected<E\&>} object does not propagate to the referenced object.
\end{itemdescr}

\indexlibrarymember{swap}{unexpected}%
\begin{itemdecl}
constexpr void swap(unexpected& other) noexcept;
\end{itemdecl}

\begin{itemdescr}
\pnum
\effects
Exchanges \exposid{unex} and \tcode{other.\exposid{unex}}.
\end{itemdescr}

\indexlibraryglobal{operator==}%
\begin{itemdecl}
template<class E2>
  friend constexpr bool operator==(const unexpected& x, const unexpected<E2>& y);
\end{itemdecl}

\begin{itemdescr}
\pnum
\mandates
The expression \tcode{x.error() == y.error()} is well-formed and its result is convertible to \tcode{bool}.

\pnum
\returns
\tcode{x.error() == y.error()}.
\end{itemdescr}

\indexlibraryglobal{swap}%
\begin{itemdecl}
friend constexpr void swap(unexpected& x, unexpected& y) noexcept;
\end{itemdecl}

\begin{itemdescr}
\pnum
\effects
Equivalent to \tcode{x.swap(y)}.
\end{itemdescr}

\rSec2[expected.bad]{Class template bad_expected_access}

\indexlibrarymember{what}{bad_expected_access}%
\begin{itemdecl}
const char* what() const noexcept override;
\end{itemdecl}

\begin{itemdescr}
\pnum
\returns
An implementation-defined ntbs, which during constant evaluation is encoded with the ordinary literal encoding (\iref{lex.ccon}).
\end{itemdescr}

\begin{codeblock}
namespace std {
  template<class E>
  class @\libglobal{bad_expected_access}@ : public bad_expected_access<void> {
  public:
    BEMAN_EXPECTED_CONSTEXPR_EXCEPTION explicit bad_expected_access(E e);
    BEMAN_EXPECTED_CONSTEXPR_EXCEPTION const char* what() const noexcept override;
    constexpr E& error() & noexcept;
    constexpr const E& error() const& noexcept;
    constexpr E&& error() && noexcept;
    constexpr const E&& error() const&& noexcept;

  private:
    E @\exposid{unex}@; // \expos
  };
}
\end{codeblock}

\pnum
\remarks
The class template \tcode{bad_expected_access} defines the type of objects thrown as exceptions to report the situation where an attempt is made to access the value of an \tcode{expected<T, E>} object for which \tcode{has_value()} is \tcode{false}.

\indexlibraryctor{bad_expected_access}%
\begin{itemdecl}
explicit bad_expected_access(E e);
\end{itemdecl}

\begin{itemdescr}
\pnum
\effects
Initializes \exposid{unex} with \tcode{std::move(e)}.
\end{itemdescr}

\indexlibrarymember{what}{bad_expected_access}%
\begin{itemdecl}
const char* what() const noexcept override;
\end{itemdecl}

\begin{itemdescr}
\pnum
\returns
An implementation-defined ntbs, which during constant evaluation is encoded with the ordinary literal encoding (\iref{lex.ccon}).
\end{itemdescr}

\indexlibrarymember{error}{bad_expected_access}%
\begin{itemdecl}
constexpr E& error() & noexcept;
\end{itemdecl}

\begin{itemdescr}
\pnum
\returns
\exposid{unex}.
\end{itemdescr}

\indexlibrarymember{error}{bad_expected_access}%
\begin{itemdecl}
constexpr const E& error() const& noexcept;
\end{itemdecl}

\begin{itemdescr}
\pnum
\returns
\exposid{unex}.
\end{itemdescr}

\indexlibrarymember{error}{bad_expected_access}%
\begin{itemdecl}
constexpr E&& error() && noexcept;
\end{itemdecl}

\begin{itemdescr}
\pnum
\returns
\tcode{std::move(\exposid{unex})}.
\end{itemdescr}

\indexlibrarymember{error}{bad_expected_access}%
\begin{itemdecl}
constexpr const E&& error() const&& noexcept;
\end{itemdecl}

\begin{itemdescr}
\pnum
\returns
\tcode{std::move(\exposid{unex})}.
\end{itemdescr}

\rSec2[expected.bad.void]{Class template specialization bad_expected_access<void>}

\begin{codeblock}
namespace std {
  template<>
  class @\libglobal{bad_expected_access}@<void> : public exception {
  protected:
    BEMAN_EXPECTED_CONSTEXPR_EXCEPTION bad_expected_access() noexcept = default;
    BEMAN_EXPECTED_CONSTEXPR_EXCEPTION bad_expected_access(
        const bad_expected_access&) noexcept = default;
    BEMAN_EXPECTED_CONSTEXPR_EXCEPTION bad_expected_access(
        bad_expected_access&&) noexcept = default;
    BEMAN_EXPECTED_CONSTEXPR_EXCEPTION bad_expected_access& operator=(
        const bad_expected_access&) noexcept = default;
    BEMAN_EXPECTED_CONSTEXPR_EXCEPTION bad_expected_access& operator=(
        bad_expected_access&&) noexcept = default;
    BEMAN_EXPECTED_CONSTEXPR_EXCEPTION ~bad_expected_access() = default;

  public:
    BEMAN_EXPECTED_CONSTEXPR_EXCEPTION const char* what() const noexcept override;
  };
}
\end{codeblock}

\rSec2[expected.expected]{Class template expected}

\rSec3[expected.object.general]{General}

\begin{codeblock}
namespace std {
  template<class T, class E>
  class @\libglobal{expected}@ {
  private:
    using @\exposid{error-value-type}@ = remove_cv_t<remove_reference_t<E>>; // \expos

  public:
    using @\libmember{value_type}{expected}@ = T;
    using @\libmember{error_type}{expected}@ = E;
    using @\libmember{unexpected_type}{expected}@ = unexpected<E>;

    template<class U> using rebind = expected<U, error_type>;

    constexpr expected() noexcept(is_nothrow_default_constructible_v<T>);

    constexpr expected(const expected&) = default;

    constexpr expected(const expected& rhs) noexcept(
        is_nothrow_copy_constructible_v<T> && is_nothrow_copy_constructible_v<E>);

    constexpr expected(expected&&) = default;

    constexpr expected(expected&& rhs) noexcept(is_nothrow_move_constructible_v<T> &&
                                                is_nothrow_move_constructible_v<E>);

    template<class U, class G>
    constexpr explicit(!is_convertible_v<const U&, T> || !is_convertible_v<const G&, E>)
        expected(const expected<U, G>& rhs);

    template<class U, class G>
    constexpr explicit(!is_convertible_v<U, T> || !is_convertible_v<G, E>)
        expected(expected<U, G>&& rhs);

    template<class U, class G>
    constexpr explicit(!is_convertible_v<const U&, T> || !is_convertible_v<G, E>)
        expected(const expected<U, G>& rhs);

    template<class U, class G>
    constexpr explicit(!is_convertible_v<U&&, T> || !is_convertible_v<G, E>)
        expected(expected<U, G>&& rhs);

    template<class U = remove_cv_t<T>>
    constexpr explicit(!is_convertible_v<U, T>) expected(U&& v);

    template<class G>
    constexpr explicit(!is_convertible_v<const G&, E>) expected(const unexpected<G>& e);

    template<class G>
    constexpr explicit(!is_convertible_v<G, E>) expected(unexpected<G>&& e);

    template<class G>
    constexpr explicit(!is_convertible_v<G, E>)
        expected(const unexpected<G>& e) noexcept;

    template<class G>
    constexpr explicit(!is_convertible_v<G, E>) expected(unexpected<G>&& e) noexcept;

    template<class G> constexpr expected(const unexpected<G>&) = delete;

    template<class G>
      requires(is_reference_v<E> && !is_reference_v<G>)
    constexpr expected(unexpected<G>&&) = delete;

    template<class... Args> constexpr explicit expected(in_place_t, Args&&... args);

    template<class U, class... Args>
    constexpr explicit expected(in_place_t, initializer_list<U> il, Args&&... args);

    template<class... Args> constexpr explicit expected(unexpect_t, Args&&... args);

    template<class... Args> constexpr expected(unexpect_t, Args&&...) = delete;

    template<class... Args>
      requires(is_reference_v<E> && !is_constructible_v<E, Args...> &&
               !@\exposid{unexpect-dangles-v}@<E, Args...>)
    constexpr expected(unexpect_t, Args&&...) = delete;

    template<class U, class... Args>
    constexpr explicit expected(unexpect_t, initializer_list<U> il, Args&&... args);

    template<class U, class... Args>
    constexpr expected(unexpect_t, initializer_list<U>, Args&&...) = delete;

    constexpr ~expected() = default;

    constexpr ~expected();

    constexpr expected& operator=(const expected&) = default;

    constexpr expected& operator=(const expected& rhs) noexcept(
        is_nothrow_copy_constructible_v<T> && is_nothrow_copy_assignable_v<T> &&
        is_nothrow_copy_constructible_v<E> && is_nothrow_copy_assignable_v<E>);

    constexpr expected& operator=(expected&&) noexcept = default;

    constexpr expected& operator=(expected&& rhs) noexcept(
        is_nothrow_move_constructible_v<T> && is_nothrow_move_assignable_v<T> &&
        is_nothrow_move_constructible_v<E> && is_nothrow_move_assignable_v<E>);

    template<class U = remove_cv_t<T>> constexpr expected& operator=(U&& v);

    template<class G> constexpr expected& operator=(const unexpected<G>& e);

    template<class G> constexpr expected& operator=(unexpected<G>&& e);

    template<class G> constexpr expected& operator=(const unexpected<G>& e);

    template<class G> constexpr expected& operator=(unexpected<G>&& e);

    template<class G> constexpr expected& operator=(const unexpected<G>&) = delete;

    template<class G>
      requires(is_reference_v<E> && !is_reference_v<G>)
    constexpr expected& operator=(unexpected<G>&&) = delete;

    template<class... Args> constexpr T& emplace(Args&&... args) noexcept;

    template<class U, class... Args>
    constexpr T& emplace(initializer_list<U> il, Args&&... args) noexcept;

    constexpr void swap(expected& rhs) noexcept(is_nothrow_move_constructible_v<T> &&
                                                is_nothrow_swappable_v<T> &&
                                                is_nothrow_move_constructible_v<E> &&
                                                (is_reference_v<E> ||
                                                 is_nothrow_swappable_v<E>));

    friend constexpr void swap(expected& x, expected& y) noexcept(noexcept(x.swap(y)));

    constexpr const T* operator->() const noexcept;
    constexpr T* operator->() noexcept;

    constexpr const T& operator*() const& noexcept;
    constexpr T& operator*() & noexcept;
    constexpr const T&& operator*() const&& noexcept;
    constexpr T&& operator*() && noexcept;

    constexpr explicit operator bool() const noexcept;
    constexpr bool has_value() const noexcept;

    constexpr const T& value() const&;
    constexpr T& value() &;
    constexpr const T&& value() const&&;
    constexpr T&& value() &&;

    constexpr const E& error() const& noexcept;
    constexpr E& error() & noexcept;
    constexpr const E&& error() const&& noexcept;
    constexpr E&& error() && noexcept;

    template<class U = remove_cv_t<T>> constexpr T value_or(U&& def) const&;

    template<class U = remove_cv_t<T>> constexpr T value_or(U&& def) &&;

    template<class G = @\exposid{error-value-type}@>
    constexpr @\exposid{error-value-type}@ error_or(G&& def) const&;

    template<class G = @\exposid{error-value-type}@>
    constexpr @\exposid{error-value-type}@ error_or(G&& def) &&;

    template<class F> constexpr auto and_then(F&& f) &;
    template<class F> constexpr auto and_then(F&& f) &&;
    template<class F> constexpr auto and_then(F&& f) const&;
    template<class F> constexpr auto and_then(F&& f) const&&;

    template<class F> constexpr auto or_else(F&& f) &;
    template<class F> constexpr auto or_else(F&& f) &&;
    template<class F> constexpr auto or_else(F&& f) const&;
    template<class F> constexpr auto or_else(F&& f) const&&;

    template<class F> constexpr auto transform(F&& f) &;
    template<class F> constexpr auto transform(F&& f) &&;
    template<class F> constexpr auto transform(F&& f) const&;
    template<class F> constexpr auto transform(F&& f) const&&;

    template<class F> constexpr auto transform_error(F&& f) &;
    template<class F> constexpr auto transform_error(F&& f) &&;
    template<class F> constexpr auto transform_error(F&& f) const&;
    template<class F> constexpr auto transform_error(F&& f) const&&;

    template<class T2, class E2>
    friend constexpr bool operator==(const expected& x, const expected<T2, E2>& y);

    template<class T2>
    friend constexpr bool operator==(const expected& x, const T2& val);

    template<class E2>
    friend constexpr bool operator==(const expected& x, const unexpected<E2>& e);

  private:
    bool @\exposid{has_val}@; // \expos
    union {
      T @\exposid{val}@;              // \expos
      unexpected<E> @\exposid{unex}@; // \expos
    };
  };
}
\end{codeblock}

\pnum
\mandates
A program that instantiates the definition of \tcode{expected<T, E>} with a \tcode{T} that is not a valid value type for \tcode{expected} (that is, \tcode{remove_cv_t<T>} is \tcode{void}, or a complete non-array object type other than \tcode{in_place_t}, \tcode{unexpect_t}, or a specialization of \tcode{unexpected}) is ill-formed. A program that instantiates the definition of \tcode{expected<T, E>} with an \tcode{E} that is not a valid template argument for \tcode{unexpected} is ill-formed.

\pnum
\remarks
Any object of type \tcode{expected<T, E>} either contains a value of type \tcode{T} or a value of type \tcode{E} nested within it. Member \exposid{has_val} indicates whether the \tcode{expected<T, E>} object contains an object of type \tcode{T}. When \tcode{has_value()} is \tcode{false}, the error is \tcode{\exposid{unex}.error()}. The error is held as an \tcode{unexpected<E>}, and not as an \tcode{E}, so that \tcode{E} may be an lvalue reference type: an \tcode{E\&} cannot be a union member, whereas \tcode{unexpected<E\&>} holds a pointer to an external object.

\rSec3[expected.object.cons]{Constructors}

\indexlibraryctor{expected}%
\begin{itemdecl}
constexpr expected() noexcept(is_nothrow_default_constructible_v<T>);
\end{itemdecl}

\begin{itemdescr}
\pnum
\constraints
\tcode{is_default_constructible_v<T>} is \tcode{true}.

\pnum
\effects
Value-initializes \exposid{val}.

\pnum
\ensures
\tcode{has_value()} is \tcode{true}.

\pnum
\throws
Any exception thrown by the initialization of \exposid{val}.
\end{itemdescr}

\indexlibraryctor{expected}%
\begin{itemdecl}
constexpr expected(const expected&) = default;
\end{itemdecl}

\begin{itemdescr}
\pnum
\effects
Direct-non-list-initializes \exposid{val} or \exposid{unex} (matching \tcode{rhs}'s active member) by trivial copy construction.

\pnum
\ensures
\tcode{rhs.has_value() == this->has_value()}.

\pnum
\remarks
This constructor is trivial.
\end{itemdescr}

\indexlibraryctor{expected}%
\begin{itemdecl}
constexpr expected(const expected& rhs) noexcept(is_nothrow_copy_constructible_v<T> &&
                                                 is_nothrow_copy_constructible_v<E>);
\end{itemdecl}

\begin{itemdescr}
\pnum
\constraints
\tcode{is_copy_constructible_v<T>} is \tcode{true}, \tcode{is_copy_constructible_v<E>} is \tcode{true}, and \tcode{(is_trivially_copy_constructible_v<T> \&\& is_trivially_copy_constructible_v<E>)} is \tcode{false}.

\pnum
\effects
If \tcode{rhs.has_value()} is \tcode{true}, direct-non-list-initializes \exposid{val} with \tcode{*rhs}. Otherwise, direct-non-list-initializes \exposid{unex} with \tcode{rhs.error()}.

\pnum
\ensures
\tcode{rhs.has_value() == this->has_value()}.

\pnum
\throws
Any exception thrown by the initialization of \exposid{val} or \exposid{unex}.

\pnum
\remarks
This constructor is defined as deleted unless \tcode{is_copy_constructible_v<T>} is \tcode{true} and \tcode{is_copy_constructible_v<E>} is \tcode{true} or \tcode{is_reference_v<E>} is \tcode{true}. This constructor is trivial if \tcode{is_trivially_copy_constructible_v<T>} is \tcode{true} and \tcode{is_trivially_copy_constructible_v<E>} is \tcode{true}.
\end{itemdescr}

\indexlibraryctor{expected}%
\begin{itemdecl}
constexpr expected(expected&&) = default;
\end{itemdecl}

\begin{itemdescr}
\pnum
\effects
Direct-non-list-initializes \exposid{val} or \exposid{unex} (matching \tcode{rhs}'s active member) by trivial move construction.

\pnum
\ensures
\tcode{rhs.has_value() == this->has_value()}.

\pnum
\remarks
This constructor is trivial.
\end{itemdescr}

\indexlibraryctor{expected}%
\begin{itemdecl}
constexpr expected(expected&& rhs) noexcept(is_nothrow_move_constructible_v<T> && is_nothrow_move_constructible_v<E>);
\end{itemdecl}

\begin{itemdescr}
\pnum
\constraints
\tcode{is_move_constructible_v<T>} is \tcode{true}, \tcode{is_move_constructible_v<E>} is \tcode{true}, and \tcode{(is_trivially_move_constructible_v<T> \&\& is_trivially_move_constructible_v<E>)} is \tcode{false}.

\pnum
\effects
If \tcode{rhs.has_value()} is \tcode{true}, direct-non-list-initializes \exposid{val} with \tcode{std::move(*rhs)}. Otherwise, direct-non-list-initializes \exposid{unex} with \tcode{std::move(rhs.error())}.

\pnum
\ensures
\tcode{rhs.has_value()} is unchanged; \tcode{rhs.has_value() == this->has_value()} is \tcode{true}.

\pnum
\throws
Any exception thrown by the initialization of \exposid{val} or \exposid{unex}.

\pnum
\remarks
This constructor is trivial if \tcode{is_trivially_move_constructible_v<T>} is \tcode{true} and \tcode{is_trivially_move_constructible_v<E>} is \tcode{true}.
\end{itemdescr}

\indexlibraryctor{expected}%
\begin{itemdecl}
template<class U, class G>
  constexpr explicit(!is_convertible_v<const U&, T> || !is_convertible_v<const G&, E>)
      expected(const expected<U, G>& rhs);
template<class U, class G>
  constexpr explicit(!is_convertible_v<U, T> || !is_convertible_v<G, E>) expected(expected<U, G>&& rhs);
\end{itemdecl}

\begin{itemdescr}
\pnum
\constraints
\tcode{is_constructible_v<T, const U\&>} is \tcode{true}; and \tcode{is_constructible_v<E, const G\&>} is \tcode{true}; and if \tcode{T} is not \tcode{bool}, \tcode{converts-from-any-cvref<T, expected<U, G>>} is \tcode{false}; and \tcode{is_constructible_v<unexpected<E>, expected<U, G>\&>} is \tcode{false}; and \tcode{is_constructible_v<unexpected<E>, expected<U, G>>} is \tcode{false}; and \tcode{is_constructible_v<unexpected<E>, const expected<U, G>\&>} is \tcode{false}; and \tcode{is_constructible_v<unexpected<E>, const expected<U, G>>} is \tcode{false}.

\pnum
\effects
If \tcode{rhs.has_value()}, direct-non-list-initializes \exposid{val} with \tcode{*rhs}. Otherwise, direct-non-list-initializes \exposid{unex} with \tcode{rhs.error()}.

\pnum
\ensures
\tcode{rhs.has_value()} is unchanged; \tcode{rhs.has_value() == this->has_value()} is \tcode{true}.

\pnum
\throws
Any exception thrown by the initialization of \exposid{val} or \exposid{unex}.
\end{itemdescr}

\indexlibraryctor{expected}%
\begin{itemdecl}
template<class U, class G>
  constexpr explicit(!is_convertible_v<const U&, T> || !is_convertible_v<G, E>) expected(const expected<U, G>& rhs);
template<class U, class G>
  constexpr explicit(!is_convertible_v<U&&, T> || !is_convertible_v<G, E>) expected(expected<U, G>&& rhs);
\end{itemdecl}

\begin{itemdescr}
\pnum
\constraints
\tcode{is_constructible_v<T, const U\&>} is \tcode{true}; and \tcode{is_reference_v<G>} is \tcode{true}; and \tcode{is_convertible_v<G, E>} is \tcode{true}.

\pnum
\effects
If \tcode{rhs.has_value()}, direct-non-list-initializes \exposid{val} with \tcode{*rhs}. Otherwise, direct-non-list-initializes \exposid{unex} with \tcode{rhs.error()}.

\pnum
\ensures
\tcode{rhs.has_value()} is unchanged; \tcode{rhs.has_value() == this->has_value()} is \tcode{true}.

\pnum
\throws
Any exception thrown by the initialization of \exposid{val} or \exposid{unex}.

\pnum
\remarks
Unlike the value-\tcode{E} overload above, this overload participates in overload resolution only when \tcode{E} and \tcode{G} are both reference types, so the referenced error object is never copied.
\end{itemdescr}

\indexlibraryctor{expected}%
\begin{itemdecl}
template<class U = remove_cv_t<T>>
  constexpr explicit(!is_convertible_v<U, T>) expected(U&& v);
\end{itemdecl}

\begin{itemdescr}
\pnum
\constraints
\tcode{is_same_v<remove_cvref_t<U>, in_place_t>} is \tcode{false}; and \tcode{is_same_v<remove_cvref_t<U>, unexpect_t>} is \tcode{false}; and \tcode{is_same_v<remove_cvref_t<U>, expected>} is \tcode{false}; and \tcode{is_constructible_v<T, U>} is \tcode{true}; and \tcode{remove_cvref_t<U>} is not a specialization of \tcode{unexpected}; and if \tcode{T} is \tcode{bool}, \tcode{remove_cvref_t<U>} is not a specialization of \tcode{expected}.

\pnum
\effects
Direct-non-list-initializes \exposid{val} with \tcode{std::forward<U>(v)}.

\pnum
\ensures
\tcode{has_value()} is \tcode{true}.

\pnum
\throws
Any exception thrown by the initialization of \exposid{val}.
\end{itemdescr}

\indexlibraryctor{expected}%
\begin{itemdecl}
template<class G>
  constexpr explicit(!is_convertible_v<const G&, E>) expected(const unexpected<G>& e);
template<class G>
  constexpr explicit(!is_convertible_v<G, E>) expected(unexpected<G>&& e);
\end{itemdecl}

\begin{itemdescr}
\pnum
\constraints
\tcode{is_constructible_v<E, const G\&>} is \tcode{true}.

\pnum
\effects
Direct-non-list-initializes \exposid{unex} with \tcode{std::forward<const G\&>(e.error())}.

\pnum
\ensures
\tcode{has_value()} is \tcode{false}.

\pnum
\throws
Any exception thrown by the initialization of \exposid{unex}.
\end{itemdescr}

\indexlibraryctor{expected}%
\begin{itemdecl}
template<class G>
  constexpr explicit(!is_convertible_v<G, E>) expected(const unexpected<G>& e) noexcept;
template<class G>
  constexpr explicit(!is_convertible_v<G, E>) expected(unexpected<G>&& e) noexcept;
\end{itemdecl}

\begin{itemdescr}
\pnum
\constraints
\tcode{is_reference_v<G>} is \tcode{true}; and \tcode{is_constructible_v<E, G>} is \tcode{true}; and \tcode{reference_constructs_from_temporary_v<E, G>} is \tcode{false}.

\pnum
\effects
Initializes \exposid{unex} with \tcode{e.error()}.

\pnum
\ensures
\tcode{has_value()} is \tcode{false}.

\pnum
\remarks
This constructor never throws: the referent is bound, not copied.
\end{itemdescr}

\indexlibraryctor{expected}%
\begin{itemdecl}
template<class G>
  constexpr expected(const unexpected<G>&) = delete;
\end{itemdecl}

\begin{itemdescr}
\pnum
\constraints
\tcode{is_reference_v<E>} is \tcode{true} and \tcode{is_reference_v<G>} is \tcode{false}.

\pnum
\remarks
When \tcode{E} is a reference type, an overload taking \tcode{unexpected<G>} for a non-reference \tcode{G} is defined as deleted: the referent would live inside the (possibly temporary) source \tcode{unexpected<G>} object, and binding \tcode{E\&} to it would dangle. Use an \tcode{unexpected<E\&>} holding an external object instead.
\end{itemdescr}

\indexlibraryctor{expected}%
\begin{itemdecl}
template<class... Args>
  constexpr explicit expected(in_place_t, Args&&... args);
\end{itemdecl}

\begin{itemdescr}
\pnum
\constraints
\tcode{is_constructible_v<T, Args...>} is \tcode{true}.

\pnum
\effects
Direct-non-list-initializes \exposid{val} with \tcode{std::forward<Args>(args)...}.

\pnum
\ensures
\tcode{has_value()} is \tcode{true}.

\pnum
\throws
Any exception thrown by the initialization of \exposid{val}.
\end{itemdescr}

\indexlibraryctor{expected}%
\begin{itemdecl}
template<class U, class... Args>
  constexpr explicit expected(in_place_t, initializer_list<U> il, Args&&... args);
\end{itemdecl}

\begin{itemdescr}
\pnum
\constraints
\tcode{is_constructible_v<T, initializer_list<U>\&, Args...>} is \tcode{true}.

\pnum
\effects
Direct-non-list-initializes \exposid{val} with \tcode{il, std::forward<Args>(args)...}.

\pnum
\ensures
\tcode{has_value()} is \tcode{true}.

\pnum
\throws
Any exception thrown by the initialization of \exposid{val}.
\end{itemdescr}

\indexlibraryctor{expected}%
\begin{itemdecl}
template<class... Args>
  constexpr explicit expected(unexpect_t, Args&&... args);
\end{itemdecl}

\begin{itemdescr}
\pnum
\constraints
\tcode{is_constructible_v<E, Args...>} is \tcode{true}.

\pnum
\effects
Direct-non-list-initializes \exposid{unex} with \tcode{std::forward<Args>(args)...}.

\pnum
\ensures
\tcode{has_value()} is \tcode{false}.

\pnum
\throws
Any exception thrown by the initialization of \exposid{unex}.
\end{itemdescr}

\indexlibraryctor{expected}%
\begin{itemdecl}
template<class... Args>
  constexpr expected(unexpect_t, Args&&...) = delete;
\end{itemdecl}

\begin{itemdescr}
\pnum
\constraints
\tcode{\exposid{unexpect-dangles-v}<E, Args...>} is \tcode{true}.

\pnum
\remarks
When \tcode{E} is a reference type, an overload with the same parameter types is defined as deleted if the single argument would bind \tcode{E\&} to a temporary, or if it is otherwise not usable to construct \tcode{E}.
\end{itemdescr}

\indexlibraryctor{expected}%
\begin{itemdecl}
template<class U, class... Args>
  constexpr explicit expected(unexpect_t, initializer_list<U> il, Args&&... args);
\end{itemdecl}

\begin{itemdescr}
\pnum
\constraints
\tcode{is_reference_v<E>} is \tcode{false}, and \tcode{is_constructible_v<E, initializer_list<U>\&, Args...>} is \tcode{true}.

\pnum
\effects
Direct-non-list-initializes \exposid{unex} with \tcode{il, std::forward<Args>(args)...}.

\pnum
\ensures
\tcode{has_value()} is \tcode{false}.

\pnum
\throws
Any exception thrown by the initialization of \exposid{unex}.

\pnum
\remarks
An overload with the same parameter types is defined as deleted when \tcode{E} is an lvalue reference type. An initializer list cannot provide the required long-lived error referent.
\end{itemdescr}

\indexlibraryctor{expected}%
\begin{itemdecl}
template<class U, class... Args>
  constexpr expected(unexpect_t, initializer_list<U>, Args&&...) = delete;
\end{itemdecl}

\begin{itemdescr}
\pnum
\constraints
\tcode{is_reference_v<E>} is \tcode{true}.

\pnum
\remarks
An overload with the same parameter types is defined as deleted when \tcode{E} is an lvalue reference type. An initializer list cannot provide the required long-lived error referent.
\end{itemdescr}

\rSec3[expected.object.dtor]{Destructor}

\indexlibrarydtor{expected}%
\begin{itemdecl}
constexpr ~expected();
\end{itemdecl}

\begin{itemdescr}
\pnum
\constraints
\tcode{is_trivially_destructible_v<T>} is \tcode{true} and \tcode{is_trivially_destructible_v<E>} is \tcode{true}.

\pnum
\effects
None: \exposid{val} or \exposid{unex} (whichever is active) has a trivial destructor.

\pnum
\remarks
This destructor is trivial.
\end{itemdescr}

\indexlibrarydtor{expected}%
\begin{itemdecl}
constexpr ~expected();
\end{itemdecl}

\begin{itemdescr}
\pnum
\constraints
\tcode{(is_trivially_destructible_v<T> \&\& is_trivially_destructible_v<E>)} is \tcode{false}.

\pnum
\effects
If \tcode{has_value()} is \tcode{true}, destroys \exposid{val}, otherwise destroys \exposid{unex}.
\end{itemdescr}

\rSec3[expected.object.assign]{Assignment}

\indexlibrarymember{operator=}{expected}%
\begin{itemdecl}
constexpr expected& operator=(const expected&);
\end{itemdecl}

\begin{itemdescr}
\pnum
\constraints
\begin{itemize}
\item \tcode{is_trivially_copy_constructible_v<T>} is \tcode{true},
\item \tcode{is_trivially_copy_assignable_v<T>} is \tcode{true},
\item \tcode{is_trivially_destructible_v<T>} is \tcode{true},
\item \tcode{is_trivially_copy_constructible_v<E>} is \tcode{true},
\item \tcode{is_trivially_copy_assignable_v<E>} is \tcode{true},
\item \tcode{is_trivially_destructible_v<E>} is \tcode{true}.
\end{itemize}

\pnum
\effects
Trivially copies \tcode{rhs}'s active member into \tcode{*this}.

\pnum
\returns
\tcode{*this}.

\pnum
\remarks
This operator is trivial.
\end{itemdescr}

\indexlibrarymember{operator=}{expected}%
\begin{itemdecl}
constexpr expected& operator=(const expected& rhs) noexcept(is_nothrow_copy_constructible_v<T> &&
                                                            is_nothrow_copy_assignable_v<T> &&
                                                            is_nothrow_copy_constructible_v<E> &&
                                                            is_nothrow_copy_assignable_v<E>);
\end{itemdecl}

\begin{itemdescr}
\pnum
\constraints
\begin{itemize}
\item \tcode{is_copy_constructible_v<T>} is \tcode{true},
\item \tcode{is_copy_assignable_v<T>} is \tcode{true},
\item \tcode{is_reference_v<E> || (is_copy_constructible_v<E> \&\& is_copy_assignable_v<E>)} is \tcode{true},
\item \tcode{is_nothrow_move_constructible_v<T> || is_nothrow_move_constructible_v<E>} is \tcode{true},
\item \tcode{(is_trivially_copy_constructible_v<T> \&\& is_trivially_copy_assignable_v<T> \&\& is_trivially_destructible_v<T> \&\& is_trivially_copy_constructible_v<E> \&\& is_trivially_copy_assignable_v<E> \&\& is_trivially_destructible_v<E>)} is \tcode{false}.
\end{itemize}

\pnum
\effects
If \tcode{this->has_value() \&\& rhs.has_value()}, equivalent to \tcode{\exposid{val} = *rhs}. Otherwise, if \tcode{this->has_value()}, equivalent to \tcode{reinit-expected(\exposid{unex}, \exposid{val}, rhs.error())}. Otherwise, if \tcode{rhs.has_value()}, equivalent to \tcode{reinit-expected(\exposid{val}, \exposid{unex}, *rhs)}. Otherwise, equivalent to \tcode{\exposid{unex} = rhs.\exposid{unex}}. Then, if no exception was thrown, equivalent to: \tcode{\exposid{has_val} = rhs.has_value(); return *this;} When \tcode{E} is an lvalue reference type, each of the cases above that initializes or assigns \exposid{unex} rebinds it: \exposid{unex} comes to refer to the same object as \tcode{rhs}'s error. No previously or subsequently referenced object is assigned through.

\pnum
\returns
\tcode{*this}.

\pnum
\remarks
This operator is defined as deleted unless \tcode{is_copy_assignable_v<T>} is \tcode{true} and \tcode{is_copy_constructible_v<T>} is \tcode{true} and \tcode{is_copy_assignable_v<E>} is \tcode{true} or \tcode{is_reference_v<E>} is \tcode{true} and \tcode{is_copy_constructible_v<E>} is \tcode{true} or \tcode{is_reference_v<E>} is \tcode{true} and \tcode{is_nothrow_move_constructible_v<T> || is_nothrow_move_constructible_v<E>} is \tcode{true}. This operator is trivial if \tcode{is_trivially_copy_constructible_v<T>}, \tcode{is_trivially_copy_assignable_v<T>}, \tcode{is_trivially_destructible_v<T>}, \tcode{is_trivially_copy_constructible_v<E>}, \tcode{is_trivially_copy_assignable_v<E>}, and \tcode{is_trivially_destructible_v<E>} are all \tcode{true}.
\end{itemdescr}

\indexlibrarymember{operator=}{expected}%
\begin{itemdecl}
constexpr expected& operator=(expected&&) noexcept;
\end{itemdecl}

\begin{itemdescr}
\pnum
\constraints
\begin{itemize}
\item \tcode{is_trivially_move_constructible_v<T>} is \tcode{true},
\item \tcode{is_trivially_move_assignable_v<T>} is \tcode{true},
\item \tcode{is_trivially_destructible_v<T>} is \tcode{true},
\item \tcode{is_trivially_move_constructible_v<E>} is \tcode{true},
\item \tcode{is_trivially_move_assignable_v<E>} is \tcode{true},
\item \tcode{is_trivially_destructible_v<E>} is \tcode{true}.
\end{itemize}

\pnum
\effects
Trivially moves \tcode{rhs}'s active member into \tcode{*this}.

\pnum
\returns
\tcode{*this}.

\pnum
\remarks
This operator is trivial.
\end{itemdescr}

\indexlibrarymember{operator=}{expected}%
\begin{itemdecl}
constexpr expected& operator=(expected&& rhs) noexcept(is_nothrow_move_constructible_v<T> &&
                                                       is_nothrow_move_assignable_v<T> &&
                                                       is_nothrow_move_constructible_v<E> &&
                                                       is_nothrow_move_assignable_v<E>);
\end{itemdecl}

\begin{itemdescr}
\pnum
\constraints
\begin{itemize}
\item \tcode{is_move_constructible_v<T>} is \tcode{true},
\item \tcode{is_move_assignable_v<T>} is \tcode{true},
\item \tcode{is_reference_v<E> || (is_move_constructible_v<E> \&\& is_move_assignable_v<E>)} is \tcode{true},
\item \tcode{is_nothrow_move_constructible_v<T> || is_nothrow_move_constructible_v<E>} is \tcode{true},
\item \tcode{(is_trivially_move_constructible_v<T> \&\& is_trivially_move_assignable_v<T> \&\& is_trivially_destructible_v<T> \&\& is_trivially_move_constructible_v<E> \&\& is_trivially_move_assignable_v<E> \&\& is_trivially_destructible_v<E>)} is \tcode{false}.
\end{itemize}

\pnum
\effects
If \tcode{this->has_value() \&\& rhs.has_value()}, equivalent to \tcode{\exposid{val} = std::move(*rhs)}. Otherwise, if \tcode{this->has_value()}, equivalent to \tcode{reinit-expected(\exposid{unex}, \exposid{val}, std::move(rhs.error()))}. Otherwise, if \tcode{rhs.has_value()}, equivalent to \tcode{reinit-expected(\exposid{val}, \exposid{unex}, std::move(*rhs))}. Otherwise, equivalent to \tcode{\exposid{unex} = std::move(rhs.\exposid{unex})}. Then, if no exception was thrown, equivalent to: \tcode{\exposid{has_val} = rhs.has_value(); return *this;} When \tcode{E} is an lvalue reference type, each of the cases above that initializes or assigns \exposid{unex} rebinds it: \exposid{unex} comes to refer to the same object as \tcode{rhs}'s error. No previously or subsequently referenced object is assigned through.

\pnum
\returns
\tcode{*this}.

\pnum
\remarks
The exception specification is equivalent to \tcode{is_nothrow_move_assignable_v<T> \&\& is_nothrow_move_constructible_v<T> \&\& is_nothrow_move_assignable_v<E> \&\& is_nothrow_move_constructible_v<E>}. This operator is trivial if \tcode{is_trivially_move_constructible_v<T>}, \tcode{is_trivially_move_assignable_v<T>}, \tcode{is_trivially_destructible_v<T>}, \tcode{is_trivially_move_constructible_v<E>}, \tcode{is_trivially_move_assignable_v<E>}, and \tcode{is_trivially_destructible_v<E>} are all \tcode{true}.
\end{itemdescr}

\indexlibrarymember{operator=}{expected}%
\begin{itemdecl}
template<class U = remove_cv_t<T>>
  constexpr expected& operator=(U&& v);
\end{itemdecl}

\begin{itemdescr}
\pnum
\constraints
\tcode{is_same_v<expected, remove_cvref_t<U>>} is \tcode{false}; and \tcode{remove_cvref_t<U>} is not a specialization of \tcode{unexpected}; and \tcode{is_constructible_v<T, U>} is \tcode{true}; and \tcode{is_assignable_v<T\&, U>} is \tcode{true}; and \tcode{is_nothrow_constructible_v<T, U> || is_nothrow_move_constructible_v<T> || is_nothrow_move_constructible_v<E>} is \tcode{true}.

\pnum
\effects
If \tcode{has_value()} is \tcode{true}, equivalent to: \tcode{\exposid{val} = std::forward<U>(v);} Otherwise, equivalent to: \tcode{reinit-expected(\exposid{val}, \exposid{unex}, std::forward<U>(v)); \exposid{has_val} = true;}

\pnum
\returns
\tcode{*this}.
\end{itemdescr}

\indexlibrarymember{operator=}{expected}%
\begin{itemdecl}
template<class G>
  constexpr expected& operator=(const unexpected<G>& e);
template<class G>
  constexpr expected& operator=(unexpected<G>&& e);
\end{itemdecl}

\begin{itemdescr}
\pnum
\constraints
\tcode{is_constructible_v<E, const G\&>} is \tcode{true}; and \tcode{is_assignable_v<E\&, const G\&>} is \tcode{true}; and \tcode{is_nothrow_constructible_v<E, const G\&> || is_nothrow_move_constructible_v<T> || is_nothrow_move_constructible_v<E>} is \tcode{true}.

\pnum
\effects
If \tcode{has_value()} is \tcode{true}, equivalent to: \tcode{reinit-expected(\exposid{unex}, \exposid{val}, e.error()); \exposid{has_val} = false;} Otherwise, equivalent to: \tcode{\exposid{unex} = unexpected<E>(e.error());}

\pnum
\returns
\tcode{*this}.
\end{itemdescr}

\indexlibrarymember{operator=}{expected}%
\begin{itemdecl}
template<class G>
  constexpr expected& operator=(const unexpected<G>& e);
template<class G>
  constexpr expected& operator=(unexpected<G>&& e);
\end{itemdecl}

\begin{itemdescr}
\pnum
\constraints
\tcode{is_reference_v<G>} is \tcode{true}; and \tcode{is_constructible_v<E, G>} is \tcode{true}; and \tcode{reference_constructs_from_temporary_v<E, G>} is \tcode{false}.

\pnum
\effects
Rebinds \exposid{unex} to refer to the same object as \tcode{e.error()}, destroying \exposid{val} first if \tcode{has_value()} is \tcode{true}.

\pnum
\ensures
\tcode{has_value()} is \tcode{false}.

\pnum
\returns
\tcode{*this}.

\pnum
\remarks
This operator never throws: the referent is bound, not copied.
\end{itemdescr}

\indexlibrarymember{operator=}{expected}%
\begin{itemdecl}
template<class G>
  constexpr expected& operator=(const unexpected<G>&) = delete;
\end{itemdecl}

\begin{itemdescr}
\pnum
\constraints
\tcode{is_reference_v<E>} is \tcode{true} and \tcode{is_reference_v<G>} is \tcode{false}.

\pnum
\remarks
When \tcode{E} is a reference type, an overload taking \tcode{unexpected<G>} for a non-reference \tcode{G} is defined as deleted: it would rebind \exposid{unex} to \tcode{unexpected<G>}'s temporary storage.
\end{itemdescr}

\indexlibrarymember{emplace}{expected}%
\begin{itemdecl}
template<class... Args>
  constexpr T& emplace(Args&&... args) noexcept;
\end{itemdecl}

\begin{itemdescr}
\pnum
\constraints
\tcode{is_nothrow_constructible_v<T, Args...>} is \tcode{true}.

\pnum
\effects
Equivalent to: \tcode{if (has_value()) \{ destroy_at(addressof(\exposid{val})); \} else \{ destroy_at(addressof(\exposid{unex})); \exposid{has_val} = true; \} return *construct_at(addressof(\exposid{val}), std::forward<Args>(args)...);}
\end{itemdescr}

\indexlibrarymember{emplace}{expected}%
\begin{itemdecl}
template<class U, class... Args>
  constexpr T& emplace(initializer_list<U> il, Args&&... args) noexcept;
\end{itemdecl}

\begin{itemdescr}
\pnum
\constraints
\tcode{is_nothrow_constructible_v<T, initializer_list<U>\&, Args...>} is \tcode{true}.

\pnum
\effects
Equivalent to: \tcode{if (has_value()) \{ destroy_at(addressof(\exposid{val})); \} else \{ destroy_at(addressof(\exposid{unex})); \exposid{has_val} = true; \} return *construct_at(addressof(\exposid{val}), il, std::forward<Args>(args)...);}
\end{itemdescr}

\rSec3[expected.object.swap]{Swap}

\indexlibrarymember{swap}{expected}%
\begin{itemdecl}
constexpr void swap(expected& rhs) noexcept(is_nothrow_move_constructible_v<T> && is_nothrow_swappable_v<T> &&
                                            is_nothrow_move_constructible_v<E> &&
                                            (is_reference_v<E> || is_nothrow_swappable_v<E>));
\end{itemdecl}

\begin{itemdescr}
\pnum
\constraints
\begin{itemize}
\item \tcode{is_swappable_v<T>} is \tcode{true},
\item \tcode{is_reference_v<E> || is_swappable_v<E>} is \tcode{true},
\item \tcode{is_move_constructible_v<T>} is \tcode{true},
\item \tcode{is_move_constructible_v<E>} is \tcode{true},
\item \tcode{is_nothrow_move_constructible_v<T> || is_nothrow_move_constructible_v<E>} is \tcode{true}.
\end{itemize}

\pnum
\effects
If \tcode{this->has_value()} and \tcode{rhs.has_value()}, equivalent to \tcode{using std::swap; swap(\exposid{val}, rhs.\exposid{val});}. If neither \tcode{*this} nor \tcode{rhs} contains a value, equivalent to \tcode{using std::swap; swap(\exposid{unex}, rhs.\exposid{unex});}. If \tcode{rhs.has_value()} is \tcode{false} and \tcode{this->has_value()} is \tcode{true}, exchanges the value and error between \tcode{*this} and \tcode{rhs} (moving through a temporary so a failed move leaves both objects unchanged), leaving \tcode{has_value()} \tcode{false} and \tcode{rhs.has_value()} \tcode{true}. If \tcode{rhs.has_value()} is \tcode{true} and \tcode{this->has_value()} is \tcode{false}, equivalent to \tcode{rhs.swap(*this)}.

\pnum
\throws
Any exception thrown by the expressions in the Effects.

\pnum
\remarks
The exception specification is equivalent to \tcode{is_nothrow_move_constructible_v<T> \&\& is_nothrow_swappable_v<T> \&\& is_nothrow_move_constructible_v<E> \&\& (is_reference_v<E> || is_nothrow_swappable_v<E>)}.
\end{itemdescr}

\indexlibraryglobal{swap}%
\begin{itemdecl}
friend constexpr void swap(expected& x, expected& y) noexcept(noexcept(x.swap(y)));
\end{itemdecl}

\begin{itemdescr}
\pnum
\constraints
\begin{itemize}
\item \tcode{is_swappable_v<T>} is \tcode{true},
\item \tcode{is_reference_v<E> || is_swappable_v<E>} is \tcode{true},
\item \tcode{is_move_constructible_v<T>} is \tcode{true},
\item \tcode{is_move_constructible_v<E>} is \tcode{true},
\item \tcode{is_nothrow_move_constructible_v<T> || is_nothrow_move_constructible_v<E>} is \tcode{true}.
\end{itemize}

\pnum
\effects
Equivalent to \tcode{x.swap(y)}.
\end{itemdescr}

\rSec3[expected.object.obs]{Observers}

\indexlibrarymember{operator->}{expected}%
\begin{itemdecl}
constexpr const T* operator->() const noexcept;
constexpr T* operator->() noexcept;
\end{itemdecl}

\begin{itemdescr}
\pnum
\hardexpects
\tcode{has_value()} is \tcode{true}.

\pnum
\returns
\tcode{addressof(\exposid{val})}.
\end{itemdescr}

\indexlibrarymember{operator*}{expected}%
\begin{itemdecl}
constexpr const T& operator*() const& noexcept;
constexpr T& operator*() & noexcept;
\end{itemdecl}

\begin{itemdescr}
\pnum
\hardexpects
\tcode{has_value()} is \tcode{true}.

\pnum
\returns
\exposid{val}.
\end{itemdescr}

\indexlibrarymember{operator*}{expected}%
\begin{itemdecl}
constexpr const T&& operator*() const&& noexcept;
constexpr T&& operator*() && noexcept;
\end{itemdecl}

\begin{itemdescr}
\pnum
\hardexpects
\tcode{has_value()} is \tcode{true}.

\pnum
\returns
\tcode{std::move(\exposid{val})}.
\end{itemdescr}

\indexlibrarymember{operator bool}{expected}%
\indexlibrarymember{has_value}{expected}%
\begin{itemdecl}
constexpr explicit operator bool() const noexcept;
constexpr bool has_value() const noexcept;
\end{itemdecl}

\begin{itemdescr}
\pnum
\returns
\exposid{has_val}.
\end{itemdescr}

\indexlibrarymember{value}{expected}%
\begin{itemdecl}
constexpr const T& value() const&;
constexpr T& value() &;
\end{itemdecl}

\begin{itemdescr}
\pnum
\mandates
\tcode{is_copy_constructible_v<E>} is \tcode{true}.

\pnum
\returns
\exposid{val}, if \tcode{has_value()} is \tcode{true}.

\pnum
\throws
\tcode{bad_expected_access(as_const(error()))} if \tcode{has_value()} is \tcode{false}.
\end{itemdescr}

\indexlibrarymember{value}{expected}%
\begin{itemdecl}
constexpr const T&& value() const&&;
constexpr T&& value() &&;
\end{itemdecl}

\begin{itemdescr}
\pnum
\mandates
\tcode{is_copy_constructible_v<E>} is \tcode{true} and \tcode{is_constructible_v<E, decltype(std::move(error()))>} is \tcode{true}.

\pnum
\returns
\tcode{std::move(\exposid{val})}, if \tcode{has_value()} is \tcode{true}.

\pnum
\throws
\tcode{bad_expected_access(std::move(error()))} if \tcode{has_value()} is \tcode{false}.
\end{itemdescr}

\indexlibrarymember{error}{expected}%
\begin{itemdecl}
constexpr const E& error() const& noexcept;
constexpr E& error() & noexcept;
\end{itemdecl}

\begin{itemdescr}
\pnum
\hardexpects
\tcode{has_value()} is \tcode{false}.

\pnum
\returns
\tcode{\exposid{unex}.error()}.
\end{itemdescr}

\indexlibrarymember{error}{expected}%
\begin{itemdecl}
constexpr const E&& error() const&& noexcept;
constexpr E&& error() && noexcept;
\end{itemdecl}

\begin{itemdescr}
\pnum
\hardexpects
\tcode{has_value()} is \tcode{false}.

\pnum
\returns
\tcode{std::move(\exposid{unex}).error()}.
\end{itemdescr}

\indexlibrarymember{value_or}{expected}%
\begin{itemdecl}
template<class U = remove_cv_t<T>>
  constexpr T value_or(U&& def) const&;
\end{itemdecl}

\begin{itemdescr}
\pnum
\mandates
\tcode{is_copy_constructible_v<T>} is \tcode{true} and \tcode{is_convertible_v<U, T>} is \tcode{true}.

\pnum
\returns
\tcode{has_value() ? **this : static_cast<T>(std::forward<U>(def))}.
\end{itemdescr}

\indexlibrarymember{value_or}{expected}%
\begin{itemdecl}
template<class U = remove_cv_t<T>>
  constexpr T value_or(U&& def) &&;
\end{itemdecl}

\begin{itemdescr}
\pnum
\mandates
\tcode{is_move_constructible_v<T>} is \tcode{true} and \tcode{is_convertible_v<U, T>} is \tcode{true}.

\pnum
\returns
\tcode{has_value() ? std::move(**this) : static_cast<T>(std::forward<U>(def))}.
\end{itemdescr}

\indexlibrarymember{error_or}{expected}%
\begin{itemdecl}
template<class G = @\exposid{error-value-type}@>
  constexpr @\exposid{error-value-type}@ error_or(G&& def) const&;
\end{itemdecl}

\begin{itemdescr}
\pnum
\constraints
\tcode{is_copy_constructible_v<typename expected<T, E>::\exposid{error-value-type}>} is \tcode{true} and \tcode{is_convertible_v<G, typename expected<T, E>::\exposid{error-value-type}>} is \tcode{true}.

\pnum
\mandates
\tcode{is_copy_constructible_v<\exposid{error-value-type}>} is \tcode{true} and \tcode{is_convertible_v<G, \exposid{error-value-type}>} is \tcode{true}.

\pnum
\returns
\tcode{std::forward<G>(def)} if \tcode{has_value()} is \tcode{true}, \tcode{error()} otherwise.
\end{itemdescr}

\indexlibrarymember{error_or}{expected}%
\begin{itemdecl}
template<class G = @\exposid{error-value-type}@>
  constexpr @\exposid{error-value-type}@ error_or(G&& def) &&;
\end{itemdecl}

\begin{itemdescr}
\pnum
\constraints
\tcode{is_move_constructible_v<typename expected<T, E>::\exposid{error-value-type}>} is \tcode{true} and \tcode{is_convertible_v<G, typename expected<T, E>::\exposid{error-value-type}>} is \tcode{true}.

\pnum
\mandates
\tcode{is_move_constructible_v<\exposid{error-value-type}>} is \tcode{true} and \tcode{is_convertible_v<G, \exposid{error-value-type}>} is \tcode{true}.

\pnum
\returns
\tcode{std::forward<G>(def)} if \tcode{has_value()} is \tcode{true}, \tcode{std::move(error())} otherwise.
\end{itemdescr}

\rSec3[expected.object.monadic]{Monadic operations}

\indexlibrarymember{and_then}{expected}%
\begin{itemdecl}
template<class F>
  constexpr auto and_then(F&& f) &;
template<class F>
  constexpr auto and_then(F&& f) const&;
\end{itemdecl}

\begin{itemdescr}
\pnum
\constraints
\tcode{is_constructible_v<E, decltype(error())>} is \tcode{true}.

\pnum
\mandates
\tcode{remove_cvref_t<invoke_result_t<F, decltype((\exposid{val}))>>} is a specialization of \tcode{expected} and its \tcode{error_type} is the same type as \tcode{E}.

\pnum
\effects
Equivalent to: \tcode{if (has_value()) return invoke(std::forward<F>(f), \exposid{val}); else return U(unexpect, error());} where \tcode{U} is \tcode{remove_cvref_t<invoke_result_t<F, decltype((\exposid{val}))>>}.
\end{itemdescr}

\indexlibrarymember{and_then}{expected}%
\begin{itemdecl}
template<class F>
  constexpr auto and_then(F&& f) &&;
template<class F>
  constexpr auto and_then(F&& f) const&&;
\end{itemdecl}

\begin{itemdescr}
\pnum
\constraints
\tcode{is_constructible_v<E, decltype(std::move(error()))>} is \tcode{true}.

\pnum
\mandates
\tcode{remove_cvref_t<invoke_result_t<F, decltype(std::move(\exposid{val}))>>} is a specialization of \tcode{expected} and its \tcode{error_type} is the same type as \tcode{E}.

\pnum
\effects
Equivalent to: \tcode{if (has_value()) return invoke(std::forward<F>(f), std::move(\exposid{val})); else return U(unexpect, std::move(error()));} where \tcode{U} is \tcode{remove_cvref_t<invoke_result_t<F, decltype(std::move(\exposid{val}))>>}.
\end{itemdescr}

\indexlibrarymember{or_else}{expected}%
\begin{itemdecl}
template<class F>
  constexpr auto or_else(F&& f) &;
template<class F>
  constexpr auto or_else(F&& f) const&;
\end{itemdecl}

\begin{itemdescr}
\pnum
\constraints
\tcode{is_constructible_v<T, decltype((\exposid{val}))>} is \tcode{true}.

\pnum
\mandates
\tcode{remove_cvref_t<invoke_result_t<F, decltype(error())>>} is a specialization of \tcode{expected} and its \tcode{value_type} is the same type as \tcode{T}.

\pnum
\effects
Equivalent to: \tcode{if (has_value()) return G(in_place, \exposid{val}); else return invoke(std::forward<F>(f), error());} where \tcode{G} is \tcode{remove_cvref_t<invoke_result_t<F, decltype(error())>>}.
\end{itemdescr}

\indexlibrarymember{or_else}{expected}%
\begin{itemdecl}
template<class F>
  constexpr auto or_else(F&& f) &&;
template<class F>
  constexpr auto or_else(F&& f) const&&;
\end{itemdecl}

\begin{itemdescr}
\pnum
\constraints
\tcode{is_constructible_v<T, decltype(std::move(\exposid{val}))>} is \tcode{true}.

\pnum
\mandates
\tcode{remove_cvref_t<invoke_result_t<F, decltype(std::move(error()))>>} is a specialization of \tcode{expected} and its \tcode{value_type} is the same type as \tcode{T}.

\pnum
\effects
Equivalent to: \tcode{if (has_value()) return G(in_place, std::move(\exposid{val})); else return invoke(std::forward<F>(f), std::move(error()));} where \tcode{G} is \tcode{remove_cvref_t<invoke_result_t<F, decltype(std::move(error()))>>}.
\end{itemdescr}

\indexlibrarymember{transform}{expected}%
\begin{itemdecl}
template<class F>
  constexpr auto transform(F&& f) &;
template<class F>
  constexpr auto transform(F&& f) const&;
\end{itemdecl}

\begin{itemdescr}
\pnum
\constraints
\tcode{is_constructible_v<E, decltype(error())>} is \tcode{true}.

\pnum
\effects
Equivalent to: \tcode{if (!has_value()) return U(unexpect, error()); else return expected<U2, E>(in_place, invoke(std::forward<F>(f), \exposid{val}));} where \tcode{U2} is \tcode{remove_cv_t<invoke_result_t<F, decltype((\exposid{val}))>>} and \tcode{U} is \tcode{expected<U2, E>}.
\end{itemdescr}

\indexlibrarymember{transform}{expected}%
\begin{itemdecl}
template<class F>
  constexpr auto transform(F&& f) &&;
template<class F>
  constexpr auto transform(F&& f) const&&;
\end{itemdecl}

\begin{itemdescr}
\pnum
\constraints
\tcode{is_constructible_v<E, decltype(std::move(error()))>} is \tcode{true}.

\pnum
\effects
Equivalent to: \tcode{if (!has_value()) return U(unexpect, std::move(error())); else return expected<U2, E>(in_place, invoke(std::forward<F>(f), std::move(\exposid{val})));} where \tcode{U2} is \tcode{remove_cv_t<invoke_result_t<F, decltype(std::move(\exposid{val}))>>} and \tcode{U} is \tcode{expected<U2, E>}.
\end{itemdescr}

\indexlibrarymember{transform_error}{expected}%
\begin{itemdecl}
template<class F>
  constexpr auto transform_error(F&& f) &;
template<class F>
  constexpr auto transform_error(F&& f) const&;
\end{itemdecl}

\begin{itemdescr}
\pnum
\constraints
\tcode{is_constructible_v<T, decltype((\exposid{val}))>} is \tcode{true}.

\pnum
\mandates
\begin{itemize}
\item \tcode{is_object_v<G>} is \tcode{true},
\item \tcode{is_array_v<G>} is \tcode{false},
\item \tcode{is_same_v<G, remove_cv_t<G>>} is \tcode{true},
\item \tcode{\exposid{is-unexpected-specialization}<G>::value} is \tcode{false}.
\end{itemize}

\pnum
\effects
Equivalent to: \tcode{if (has_value()) return G(in_place, \exposid{val}); else return expected<T, G2>(unexpect, invoke(std::forward<F>(f), error()));} where \tcode{G2} is \tcode{remove_cv_t<invoke_result_t<F, decltype(error())>>} and \tcode{G} is \tcode{expected<T, G2>}.
\end{itemdescr}

\indexlibrarymember{transform_error}{expected}%
\begin{itemdecl}
template<class F>
  constexpr auto transform_error(F&& f) &&;
template<class F>
  constexpr auto transform_error(F&& f) const&&;
\end{itemdecl}

\begin{itemdescr}
\pnum
\constraints
\tcode{is_constructible_v<T, decltype(std::move(\exposid{val}))>} is \tcode{true}.

\pnum
\mandates
\begin{itemize}
\item \tcode{is_object_v<G>} is \tcode{true},
\item \tcode{is_array_v<G>} is \tcode{false},
\item \tcode{is_same_v<G, remove_cv_t<G>>} is \tcode{true},
\item \tcode{\exposid{is-unexpected-specialization}<G>::value} is \tcode{false}.
\end{itemize}

\pnum
\effects
Equivalent to: \tcode{if (has_value()) return G(in_place, std::move(\exposid{val})); else return expected<T, G2>(unexpect, invoke(std::forward<F>(f), std::move(error())));} where \tcode{G2} is \tcode{remove_cv_t<invoke_result_t<F, decltype(std::move(error()))>>} and \tcode{G} is \tcode{expected<T, G2>}.
\end{itemdescr}

\rSec3[expected.object.eq]{Equality operators}

\indexlibraryglobal{operator==}%
\begin{itemdecl}
template<class T2, class E2>
  friend constexpr bool operator==(const expected& x, const expected<T2, E2>& y);
\end{itemdecl}

\begin{itemdescr}
\pnum
\constraints
\tcode{is_void_v<T2>} is \tcode{false}.

\pnum
\mandates
\tcode{!is_void_v<T2>} is \tcode{true}. The expression \tcode{*x == *y} is well-formed and its result is convertible to \tcode{bool}. The expression \tcode{x.error() == y.error()} is well-formed and its result is convertible to \tcode{bool}.

\pnum
\returns
If \tcode{x.has_value() != y.has_value()}, \tcode{false}; otherwise, if \tcode{x.has_value()} is \tcode{true}, \tcode{*x == *y}; otherwise \tcode{x.error() == y.error()}.
\end{itemdescr}

\indexlibraryglobal{operator==}%
\begin{itemdecl}
template<class T2>
  friend constexpr bool operator==(const expected& x, const T2& val);
\end{itemdecl}

\begin{itemdescr}
\pnum
\constraints
\tcode{\exposid{is-expected-specialization}<T2>::value} is \tcode{false}.

\pnum
\mandates
\tcode{T2} is not a specialization of \tcode{expected}. The expression \tcode{*x == val} is well-formed and its result is convertible to \tcode{bool}.

\pnum
\returns
\tcode{x.has_value() \&\& static_cast<bool>(*x == val)}.
\end{itemdescr}

\indexlibraryglobal{operator==}%
\begin{itemdecl}
template<class E2>
  friend constexpr bool operator==(const expected& x, const unexpected<E2>& e);
\end{itemdecl}

\begin{itemdescr}
\pnum
\mandates
The expression \tcode{x.error() == e.error()} is well-formed and its result is convertible to \tcode{bool}.

\pnum
\returns
\tcode{!x.has_value() \&\& static_cast<bool>(x.error() == e.error())}.
\end{itemdescr}

\rSec2[expected.void]{Partial specialization of expected for void types}

\rSec3[expected.void.general]{General}

\begin{codeblock}
namespace std {
  template<class E>
  class @\libglobal{expected}@<void, E> {
  private:
    using @\exposid{error-value-type}@ = remove_cv_t<remove_reference_t<E>>; // \expos

  public:
    using @\libmember{value_type}{expected}@ = void;
    using @\libmember{error_type}{expected}@ = E;
    using @\libmember{unexpected_type}{expected}@ = unexpected<E>;

    template<class U> using rebind = expected<U, error_type>;

    constexpr expected() noexcept;

    constexpr expected(const expected& rhs) noexcept(
        is_nothrow_copy_constructible_v<E>);

    constexpr expected(expected&& rhs) noexcept(is_nothrow_move_constructible_v<E>);

    template<class U, class G>
    constexpr explicit(!is_convertible_v<const G&, E>)
        expected(const expected<U, G>& rhs);

    template<class U, class G>
    constexpr explicit(!is_convertible_v<G, E>) expected(expected<U, G>&& rhs);

    template<class G>
    constexpr explicit(!is_convertible_v<const G&, E>) expected(const unexpected<G>& e);

    template<class G>
    constexpr explicit(!is_convertible_v<G, E>) expected(unexpected<G>&& e);

    template<class G>
    constexpr explicit(!is_convertible_v<G, E>)
        expected(const unexpected<G>& e) noexcept;

    template<class G>
    constexpr explicit(!is_convertible_v<G, E>) expected(unexpected<G>&& e) noexcept;

    template<class G> constexpr expected(const unexpected<G>&) = delete;

    template<class G>
      requires(is_reference_v<E> && !is_reference_v<G>)
    constexpr expected(unexpected<G>&&) = delete;

    constexpr explicit expected(in_place_t) noexcept;

    template<class... Args> constexpr explicit expected(unexpect_t, Args&&... args);

    template<class... Args> constexpr expected(unexpect_t, Args&&...) = delete;

    template<class... Args>
      requires(is_reference_v<E> && !is_constructible_v<E, Args...> &&
               !@\exposid{unexpect-dangles-v}@<E, Args...>)
    constexpr expected(unexpect_t, Args&&...) = delete;

    template<class U, class... Args>
    constexpr explicit expected(unexpect_t, initializer_list<U> il, Args&&... args);

    template<class U, class... Args>
    constexpr expected(unexpect_t, initializer_list<U>, Args&&...) = delete;

    template<class G>
    constexpr explicit(!is_convertible_v<G&, E>)
        expected(const expected<void, G&>& rhs);

    template<class G>
    constexpr explicit(!is_convertible_v<G&, E>) expected(expected<void, G&>&& rhs);

    constexpr ~expected();

    constexpr expected& operator=(const expected& rhs) noexcept(
        is_nothrow_copy_constructible_v<E> && is_nothrow_copy_assignable_v<E>);

    constexpr expected& operator=(expected&& rhs) noexcept(
        is_nothrow_move_constructible_v<E> && is_nothrow_move_assignable_v<E>);

    template<class G> constexpr expected& operator=(const unexpected<G>& e);

    template<class G> constexpr expected& operator=(unexpected<G>&& e);

    template<class G> constexpr expected& operator=(const unexpected<G>& e);

    template<class G> constexpr expected& operator=(unexpected<G>&& e);

    template<class G> constexpr expected& operator=(const unexpected<G>&) = delete;

    template<class G>
      requires(is_reference_v<E> && !is_reference_v<G>)
    constexpr expected& operator=(unexpected<G>&&) = delete;

    constexpr void emplace() noexcept;

    constexpr void swap(expected& rhs) noexcept(is_nothrow_move_constructible_v<E> &&
                                                (is_reference_v<E> ||
                                                 is_nothrow_swappable_v<E>));

    friend constexpr void swap(expected& x, expected& y) noexcept(noexcept(x.swap(y)));

    constexpr explicit operator bool() const noexcept;
    constexpr bool has_value() const noexcept;

    constexpr void operator*() const noexcept;

    constexpr void value() const&;
    constexpr void value() &&;

    constexpr const E& error() const& noexcept;
    constexpr E& error() & noexcept;
    constexpr const E&& error() const&& noexcept;
    constexpr E&& error() && noexcept;

    template<class G = @\exposid{error-value-type}@>
    constexpr @\exposid{error-value-type}@ error_or(G&& def) const&;

    template<class G = @\exposid{error-value-type}@>
    constexpr @\exposid{error-value-type}@ error_or(G&& def) &&;

    template<class U> constexpr void value_or(U&&) const = delete;

    template<class F> constexpr auto and_then(F&& f) &;
    template<class F> constexpr auto and_then(F&& f) &&;
    template<class F> constexpr auto and_then(F&& f) const&;
    template<class F> constexpr auto and_then(F&& f) const&&;

    template<class F> constexpr auto or_else(F&& f) &;
    template<class F> constexpr auto or_else(F&& f) &&;
    template<class F> constexpr auto or_else(F&& f) const&;
    template<class F> constexpr auto or_else(F&& f) const&&;

    template<class F> constexpr auto transform(F&& f) &;
    template<class F> constexpr auto transform(F&& f) &&;
    template<class F> constexpr auto transform(F&& f) const&;
    template<class F> constexpr auto transform(F&& f) const&&;

    template<class F> constexpr auto transform_error(F&& f) &;
    template<class F> constexpr auto transform_error(F&& f) &&;
    template<class F> constexpr auto transform_error(F&& f) const&;
    template<class F> constexpr auto transform_error(F&& f) const&&;

    template<class T2, class E2>
    friend constexpr bool operator==(const expected& x, const expected<T2, E2>& y);

    template<class E2>
    friend constexpr bool operator==(const expected& x, const unexpected<E2>& e);

  private:
    bool @\exposid{has_val}@; // \expos
    union {
      unexpected<E> @\exposid{unex}@; // \expos
    };
  };
}
\end{codeblock}

\pnum
\mandates
A program that instantiates the definition of \tcode{expected<T, E>} with an \tcode{E} that is not a valid template argument for \tcode{unexpected} is ill-formed.

\pnum
\remarks
Any object of type \tcode{expected<T, E>} either represents a value of type \tcode{T}, or contains a value of type \tcode{E} nested within it. Member \exposid{has_val} indicates whether the \tcode{expected<T, E>} object represents a value of type \tcode{T}. When \tcode{has_value()} is \tcode{false}, the error is \tcode{\exposid{unex}.error()}.

\rSec3[expected.void.cons]{Constructors}

\indexlibraryctor{expected}%
\begin{itemdecl}
constexpr expected() noexcept;
\end{itemdecl}

\begin{itemdescr}
\pnum
\ensures
\tcode{has_value()} is \tcode{true}.
\end{itemdescr}

\indexlibraryctor{expected}%
\begin{itemdecl}
constexpr expected(const expected& rhs) noexcept(is_nothrow_copy_constructible_v<E>);
\end{itemdecl}

\begin{itemdescr}
\pnum
\constraints
\tcode{is_copy_constructible_v<E>} is \tcode{true} and \tcode{is_trivially_copy_constructible_v<E>} is \tcode{false}.

\pnum
\effects
If \tcode{rhs.has_value()} is \tcode{false}, direct-non-list-initializes \exposid{unex} with \tcode{rhs.error()}.

\pnum
\ensures
\tcode{rhs.has_value() == this->has_value()}.

\pnum
\throws
Any exception thrown by the initialization of \exposid{unex}.

\pnum
\remarks
This constructor is defined as deleted unless \tcode{is_copy_constructible_v<E>} is \tcode{true} or \tcode{is_reference_v<E>} is \tcode{true}. This constructor is trivial if \tcode{is_trivially_copy_constructible_v<E>} is \tcode{true}.
\end{itemdescr}

\indexlibraryctor{expected}%
\begin{itemdecl}
constexpr expected(expected&& rhs) noexcept(is_nothrow_move_constructible_v<E>);
\end{itemdecl}

\begin{itemdescr}
\pnum
\constraints
\tcode{is_move_constructible_v<E>} is \tcode{true} and \tcode{is_trivially_move_constructible_v<E>} is \tcode{false}.

\pnum
\effects
If \tcode{rhs.has_value()} is \tcode{false}, direct-non-list-initializes \exposid{unex} with \tcode{std::move(rhs.error())}.

\pnum
\ensures
\tcode{rhs.has_value()} is unchanged; \tcode{rhs.has_value() == this->has_value()} is \tcode{true}.

\pnum
\throws
Any exception thrown by the initialization of \exposid{unex}.

\pnum
\remarks
This constructor is trivial if \tcode{is_trivially_move_constructible_v<E>} is \tcode{true}.
\end{itemdescr}

\indexlibraryctor{expected}%
\begin{itemdecl}
template<class U, class G>
  constexpr explicit(!is_convertible_v<const G&, E>) expected(const expected<U, G>& rhs);
template<class U, class G>
  constexpr explicit(!is_convertible_v<G, E>) expected(expected<U, G>&& rhs);
\end{itemdecl}

\begin{itemdescr}
\pnum
\constraints
\tcode{is_void_v<U>} is \tcode{true}; and \tcode{is_constructible_v<E, const G\&>} is \tcode{true}; and \tcode{is_constructible_v<unexpected<E>, expected<U, G>\&>} is \tcode{false}; and \tcode{is_constructible_v<unexpected<E>, expected<U, G>>} is \tcode{false}; and \tcode{is_constructible_v<unexpected<E>, const expected<U, G>\&>} is \tcode{false}; and \tcode{is_constructible_v<unexpected<E>, const expected<U, G>>} is \tcode{false}.

\pnum
\effects
If \tcode{rhs.has_value()} is \tcode{false}, direct-non-list-initializes \exposid{unex} with \tcode{rhs.error()}.

\pnum
\ensures
\tcode{rhs.has_value()} is unchanged; \tcode{rhs.has_value() == this->has_value()} is \tcode{true}.

\pnum
\throws
Any exception thrown by the initialization of \exposid{unex}.
\end{itemdescr}

\indexlibraryctor{expected}%
\begin{itemdecl}
template<class G>
  constexpr explicit(!is_convertible_v<const G&, E>) expected(const unexpected<G>& e);
template<class G>
  constexpr explicit(!is_convertible_v<G, E>) expected(unexpected<G>&& e);
\end{itemdecl}

\begin{itemdescr}
\pnum
\constraints
\tcode{is_constructible_v<E, const G\&>} is \tcode{true}.

\pnum
\effects
Direct-non-list-initializes \exposid{unex} with \tcode{e.error()}.

\pnum
\ensures
\tcode{has_value()} is \tcode{false}.

\pnum
\throws
Any exception thrown by the initialization of \exposid{unex}.
\end{itemdescr}

\indexlibraryctor{expected}%
\begin{itemdecl}
template<class G>
  constexpr explicit(!is_convertible_v<G, E>) expected(const unexpected<G>& e) noexcept;
template<class G>
  constexpr explicit(!is_convertible_v<G, E>) expected(unexpected<G>&& e) noexcept;
\end{itemdecl}

\begin{itemdescr}
\pnum
\constraints
\tcode{is_reference_v<G>} is \tcode{true}; and \tcode{is_constructible_v<E, G>} is \tcode{true}; and \tcode{reference_constructs_from_temporary_v<E, G>} is \tcode{false}.

\pnum
\effects
Initializes \exposid{unex} with \tcode{e.error()}.

\pnum
\ensures
\tcode{has_value()} is \tcode{false}.

\pnum
\remarks
This constructor never throws: the referent is bound, not copied.
\end{itemdescr}

\indexlibraryctor{expected}%
\begin{itemdecl}
template<class G>
  constexpr expected(const unexpected<G>&) = delete;
\end{itemdecl}

\begin{itemdescr}
\pnum
\constraints
\tcode{is_reference_v<E>} is \tcode{true} and \tcode{is_reference_v<G>} is \tcode{false}.

\pnum
\remarks
When \tcode{E} is a reference type, an overload taking \tcode{unexpected<G>} for a non-reference \tcode{G} is defined as deleted: the referent would live inside the (possibly temporary) source \tcode{unexpected<G>} object, and binding \tcode{E\&} to it would dangle. Use an \tcode{unexpected<E\&>} holding an external object instead.
\end{itemdescr}

\indexlibraryctor{expected}%
\begin{itemdecl}
constexpr explicit expected(in_place_t) noexcept;
\end{itemdecl}

\begin{itemdescr}
\pnum
\ensures
\tcode{has_value()} is \tcode{true}.
\end{itemdescr}

\indexlibraryctor{expected}%
\begin{itemdecl}
template<class... Args>
  constexpr explicit expected(unexpect_t, Args&&... args);
\end{itemdecl}

\begin{itemdescr}
\pnum
\constraints
\tcode{is_constructible_v<E, Args...>} is \tcode{true}.

\pnum
\effects
Direct-non-list-initializes \exposid{unex} with \tcode{std::forward<Args>(args)...}.

\pnum
\ensures
\tcode{has_value()} is \tcode{false}.

\pnum
\throws
Any exception thrown by the initialization of \exposid{unex}.
\end{itemdescr}

\indexlibraryctor{expected}%
\begin{itemdecl}
template<class... Args>
  constexpr expected(unexpect_t, Args&&...) = delete;
\end{itemdecl}

\begin{itemdescr}
\pnum
\constraints
\tcode{\exposid{unexpect-dangles-v}<E, Args...>} is \tcode{true}.

\pnum
\remarks
When \tcode{E} is a reference type, an overload with the same parameter types is defined as deleted if the single argument would bind \tcode{E\&} to a temporary, or if it is otherwise not usable to construct \tcode{E}.
\end{itemdescr}

\indexlibraryctor{expected}%
\begin{itemdecl}
template<class U, class... Args>
  constexpr explicit expected(unexpect_t, initializer_list<U> il, Args&&... args);
\end{itemdecl}

\begin{itemdescr}
\pnum
\constraints
\tcode{is_reference_v<E>} is \tcode{false}, and \tcode{is_constructible_v<E, initializer_list<U>\&, Args...>} is \tcode{true}.

\pnum
\effects
Direct-non-list-initializes \exposid{unex} with \tcode{il, std::forward<Args>(args)...}.

\pnum
\ensures
\tcode{has_value()} is \tcode{false}.

\pnum
\throws
Any exception thrown by the initialization of \exposid{unex}.

\pnum
\remarks
An overload with the same parameter types is defined as deleted when \tcode{E} is an lvalue reference type. An initializer list cannot provide the required long-lived error referent.
\end{itemdescr}

\indexlibraryctor{expected}%
\begin{itemdecl}
template<class U, class... Args>
  constexpr expected(unexpect_t, initializer_list<U>, Args&&...) = delete;
\end{itemdecl}

\begin{itemdescr}
\pnum
\constraints
\tcode{is_reference_v<E>} is \tcode{true}.

\pnum
\remarks
An overload with the same parameter types is defined as deleted when \tcode{E} is an lvalue reference type. An initializer list cannot provide the required long-lived error referent.
\end{itemdescr}

\indexlibraryctor{expected}%
\begin{itemdecl}
template<class G>
  constexpr explicit(!is_convertible_v<G&, E>) expected(const expected<void, G&>& rhs);
template<class G>
  constexpr explicit(!is_convertible_v<G&, E>) expected(expected<void, G&>&& rhs);
\end{itemdecl}

\begin{itemdescr}
\pnum
\constraints
\tcode{is_convertible_v<G\&, E>} is \tcode{true}; and \tcode{reference_constructs_from_temporary_v<E, G\&>} is \tcode{false}.

\pnum
\effects
If \tcode{rhs.has_value()} is \tcode{false}, direct-non-list-initializes \exposid{unex} with \tcode{rhs.error()}.

\pnum
\ensures
\tcode{rhs.has_value()} is unchanged; \tcode{rhs.has_value() == this->has_value()} is \tcode{true}.

\pnum
\remarks
This constructor never throws: the referent is bound, not copied. It participates in overload resolution only when \tcode{E} is a reference type, mirroring the \tcode{unexpected<G>} reference-\tcode{E} path above.
\end{itemdescr}

\rSec3[expected.void.dtor]{Destructor}

\indexlibrarydtor{expected}%
\begin{itemdecl}
constexpr ~expected();
\end{itemdecl}

\begin{itemdescr}
\pnum
\effects
If \tcode{has_value()} is \tcode{false}, destroys \exposid{unex}.

\pnum
\remarks
If \tcode{is_trivially_destructible_v<E>} is \tcode{true}, then this destructor is a trivial destructor.
\end{itemdescr}

\rSec3[expected.void.assign]{Assignment}

\indexlibrarymember{operator=}{expected}%
\begin{itemdecl}
constexpr expected& operator=(const expected& rhs) noexcept(is_nothrow_copy_constructible_v<E> &&
                                                            is_nothrow_copy_assignable_v<E>);
\end{itemdecl}

\begin{itemdescr}
\pnum
\constraints
\tcode{is_reference_v<E> || (is_copy_constructible_v<E> \&\& is_copy_assignable_v<E>)} is \tcode{true} and \tcode{(is_trivially_copy_constructible_v<E> \&\& is_trivially_copy_assignable_v<E> \&\& is_trivially_destructible_v<E>)} is \tcode{false}.

\pnum
\effects
If \tcode{this->has_value() \&\& rhs.has_value()}, no effects. Otherwise, if \tcode{this->has_value()}, equivalent to: \tcode{construct_at(addressof(\exposid{unex}), rhs.\exposid{unex}); \exposid{has_val} = false;} Otherwise, if \tcode{rhs.has_value()}, destroys \exposid{unex} and sets \exposid{has_val} to \tcode{true}. Otherwise, equivalent to \tcode{\exposid{unex} = rhs.\exposid{unex}}.

\pnum
\returns
\tcode{*this}.

\pnum
\remarks
This operator is defined as deleted unless \tcode{is_copy_assignable_v<E>} is \tcode{true} or \tcode{is_reference_v<E>} is \tcode{true} and \tcode{is_copy_constructible_v<E>} is \tcode{true} or \tcode{is_reference_v<E>} is \tcode{true}. This operator is trivial if \tcode{is_trivially_copy_constructible_v<E>}, \tcode{is_trivially_copy_assignable_v<E>}, and \tcode{is_trivially_destructible_v<E>} are all \tcode{true}.
\end{itemdescr}

\indexlibrarymember{operator=}{expected}%
\begin{itemdecl}
constexpr expected& operator=(expected&& rhs) noexcept(is_nothrow_move_constructible_v<E> &&
                                                       is_nothrow_move_assignable_v<E>);
\end{itemdecl}

\begin{itemdescr}
\pnum
\constraints
\tcode{is_reference_v<E> || (is_move_constructible_v<E> \&\& is_move_assignable_v<E>)} is \tcode{true} and \tcode{(is_trivially_move_constructible_v<E> \&\& is_trivially_move_assignable_v<E> \&\& is_trivially_destructible_v<E>)} is \tcode{false}.

\pnum
\effects
If \tcode{this->has_value() \&\& rhs.has_value()}, no effects. Otherwise, if \tcode{this->has_value()}, equivalent to: \tcode{construct_at(addressof(\exposid{unex}), std::move(rhs.\exposid{unex})); \exposid{has_val} = false;} Otherwise, if \tcode{rhs.has_value()}, destroys \exposid{unex} and sets \exposid{has_val} to \tcode{true}. Otherwise, equivalent to \tcode{\exposid{unex} = std::move(rhs.\exposid{unex})}.

\pnum
\returns
\tcode{*this}.

\pnum
\remarks
The exception specification is equivalent to \tcode{is_nothrow_move_constructible_v<E> \&\& is_nothrow_move_assignable_v<E>}. This operator is trivial if \tcode{is_trivially_move_constructible_v<E>}, \tcode{is_trivially_move_assignable_v<E>}, and \tcode{is_trivially_destructible_v<E>} are all \tcode{true}.
\end{itemdescr}

\indexlibrarymember{operator=}{expected}%
\begin{itemdecl}
template<class G>
  constexpr expected& operator=(const unexpected<G>& e);
template<class G>
  constexpr expected& operator=(unexpected<G>&& e);
\end{itemdecl}

\begin{itemdescr}
\pnum
\constraints
\tcode{is_constructible_v<E, const G\&>} is \tcode{true} and \tcode{is_assignable_v<E\&, const G\&>} is \tcode{true}.

\pnum
\effects
If \tcode{has_value()} is \tcode{true}, equivalent to: \tcode{construct_at(addressof(\exposid{unex}), e.error()); \exposid{has_val} = false;} Otherwise, equivalent to: \tcode{\exposid{unex} = unexpected<E>(e.error());}

\pnum
\returns
\tcode{*this}.
\end{itemdescr}

\indexlibrarymember{operator=}{expected}%
\begin{itemdecl}
template<class G>
  constexpr expected& operator=(const unexpected<G>& e);
template<class G>
  constexpr expected& operator=(unexpected<G>&& e);
\end{itemdecl}

\begin{itemdescr}
\pnum
\constraints
\tcode{is_reference_v<G>} is \tcode{true}; and \tcode{is_constructible_v<E, G>} is \tcode{true}; and \tcode{reference_constructs_from_temporary_v<E, G>} is \tcode{false}.

\pnum
\effects
Rebinds \exposid{unex} to refer to the same object as \tcode{e.error()}.

\pnum
\ensures
\tcode{has_value()} is \tcode{false}.

\pnum
\returns
\tcode{*this}.

\pnum
\remarks
This operator never throws: the referent is bound, not copied.
\end{itemdescr}

\indexlibrarymember{operator=}{expected}%
\begin{itemdecl}
template<class G>
  constexpr expected& operator=(const unexpected<G>&) = delete;
\end{itemdecl}

\begin{itemdescr}
\pnum
\constraints
\tcode{is_reference_v<E>} is \tcode{true} and \tcode{is_reference_v<G>} is \tcode{false}.

\pnum
\remarks
When \tcode{E} is a reference type, an overload taking \tcode{unexpected<G>} for a non-reference \tcode{G} is defined as deleted: it would rebind \exposid{unex} to \tcode{unexpected<G>}'s temporary storage.
\end{itemdescr}

\indexlibrarymember{emplace}{expected}%
\begin{itemdecl}
constexpr void emplace() noexcept;
\end{itemdecl}

\begin{itemdescr}
\pnum
\effects
If \tcode{has_value()} is \tcode{false}, destroys \exposid{unex} and sets \exposid{has_val} to \tcode{true}.
\end{itemdescr}

\rSec3[expected.void.swap]{Swap}

\indexlibrarymember{swap}{expected}%
\begin{itemdecl}
constexpr void swap(expected& rhs) noexcept(is_nothrow_move_constructible_v<E> &&
                                            (is_reference_v<E> || is_nothrow_swappable_v<E>));
\end{itemdecl}

\begin{itemdescr}
\pnum
\constraints
\tcode{is_reference_v<E> || is_swappable_v<E>} is \tcode{true} and \tcode{is_move_constructible_v<E>} is \tcode{true}.

\pnum
\effects
If \tcode{this->has_value()} and \tcode{rhs.has_value()}, no effects. If neither \tcode{*this} nor \tcode{rhs} contains a value, equivalent to \tcode{using std::swap; swap(\exposid{unex}, rhs.\exposid{unex});}. If \tcode{rhs.has_value()} is \tcode{false} and \tcode{this->has_value()} is \tcode{true}, initializes \tcode{rhs.\exposid{unex}} from \tcode{std::move(\exposid{unex})}, destroys \exposid{unex}, and leaves \tcode{has_value()} \tcode{false} and \tcode{rhs.has_value()} \tcode{true}. If \tcode{rhs.has_value()} is \tcode{true} and \tcode{this->has_value()} is \tcode{false}, equivalent to \tcode{rhs.swap(*this)}.

\pnum
\throws
Any exception thrown by the expressions in the Effects.

\pnum
\remarks
The exception specification is equivalent to \tcode{is_nothrow_move_constructible_v<E> \&\& (is_reference_v<E> || is_nothrow_swappable_v<E>)}.
\end{itemdescr}

\indexlibraryglobal{swap}%
\begin{itemdecl}
friend constexpr void swap(expected& x, expected& y) noexcept(noexcept(x.swap(y)));
\end{itemdecl}

\begin{itemdescr}
\pnum
\constraints
\tcode{is_reference_v<E> || is_swappable_v<E>} is \tcode{true} and \tcode{is_move_constructible_v<E>} is \tcode{true}.

\pnum
\effects
Equivalent to \tcode{x.swap(y)}.
\end{itemdescr}

\rSec3[expected.void.obs]{Observers}

\indexlibrarymember{operator bool}{expected}%
\indexlibrarymember{has_value}{expected}%
\begin{itemdecl}
constexpr explicit operator bool() const noexcept;
constexpr bool has_value() const noexcept;
\end{itemdecl}

\begin{itemdescr}
\pnum
\returns
\exposid{has_val}.
\end{itemdescr}

\indexlibrarymember{operator*}{expected}%
\begin{itemdecl}
constexpr void operator*() const noexcept;
\end{itemdecl}

\begin{itemdescr}
\pnum
\hardexpects
\tcode{has_value()} is \tcode{true}.
\end{itemdescr}

\indexlibrarymember{value}{expected}%
\begin{itemdecl}
constexpr void value() const&;
\end{itemdecl}

\begin{itemdescr}
\pnum
\mandates
\tcode{is_copy_constructible_v<E>} is \tcode{true}.

\pnum
\throws
\tcode{bad_expected_access(error())} if \tcode{has_value()} is \tcode{false}.
\end{itemdescr}

\indexlibrarymember{value}{expected}%
\begin{itemdecl}
constexpr void value() &&;
\end{itemdecl}

\begin{itemdescr}
\pnum
\mandates
\tcode{is_copy_constructible_v<E>} is \tcode{true} and \tcode{is_move_constructible_v<E>} is \tcode{true}.

\pnum
\throws
\tcode{bad_expected_access(std::move(error()))} if \tcode{has_value()} is \tcode{false}.
\end{itemdescr}

\indexlibrarymember{error}{expected}%
\begin{itemdecl}
constexpr const E& error() const& noexcept;
constexpr E& error() & noexcept;
\end{itemdecl}

\begin{itemdescr}
\pnum
\hardexpects
\tcode{has_value()} is \tcode{false}.

\pnum
\returns
\tcode{\exposid{unex}.error()}.
\end{itemdescr}

\indexlibrarymember{error}{expected}%
\begin{itemdecl}
constexpr const E&& error() const&& noexcept;
constexpr E&& error() && noexcept;
\end{itemdecl}

\begin{itemdescr}
\pnum
\hardexpects
\tcode{has_value()} is \tcode{false}.

\pnum
\returns
\tcode{std::move(\exposid{unex}).error()}.
\end{itemdescr}

\indexlibrarymember{error_or}{expected}%
\begin{itemdecl}
template<class G = @\exposid{error-value-type}@>
  constexpr @\exposid{error-value-type}@ error_or(G&& def) const&;
\end{itemdecl}

\begin{itemdescr}
\pnum
\constraints
\tcode{is_copy_constructible_v<remove_cv_t<remove_reference_t<E>>>} is \tcode{true} and \tcode{is_convertible_v<G, remove_cv_t<remove_reference_t<E>>>} is \tcode{true}.

\pnum
\mandates
\tcode{is_copy_constructible_v<\exposid{error-value-type}>} is \tcode{true} and \tcode{is_convertible_v<G, \exposid{error-value-type}>} is \tcode{true}.

\pnum
\returns
\tcode{std::forward<G>(def)} if \tcode{has_value()} is \tcode{true}, \tcode{error()} otherwise.
\end{itemdescr}

\indexlibrarymember{error_or}{expected}%
\begin{itemdecl}
template<class G = @\exposid{error-value-type}@>
  constexpr @\exposid{error-value-type}@ error_or(G&& def) &&;
\end{itemdecl}

\begin{itemdescr}
\pnum
\constraints
\tcode{is_move_constructible_v<remove_cv_t<remove_reference_t<E>>>} is \tcode{true} and \tcode{is_convertible_v<G, remove_cv_t<remove_reference_t<E>>>} is \tcode{true}.

\pnum
\mandates
\tcode{is_move_constructible_v<\exposid{error-value-type}>} is \tcode{true} and \tcode{is_convertible_v<G, \exposid{error-value-type}>} is \tcode{true}.

\pnum
\returns
\tcode{std::forward<G>(def)} if \tcode{has_value()} is \tcode{true}, \tcode{std::move(error())} otherwise.
\end{itemdescr}

\indexlibrarymember{value_or}{expected}%
\begin{itemdecl}
template<class U>
  constexpr void value_or(U&&) const = delete;
\end{itemdecl}

\begin{itemdescr}
\pnum
\constraints
\tcode{is_reference_v<E>} is \tcode{true}.

\pnum
\remarks
\tcode{expected<void, E>} has no \tcode{value_or} member: there is no value to fall back from. This overload exists only to give a clear diagnostic when \tcode{E} is a reference type, and is defined as deleted.
\end{itemdescr}

\rSec3[expected.void.monadic]{Monadic operations}

\indexlibrarymember{and_then}{expected}%
\begin{itemdecl}
template<class F>
  constexpr auto and_then(F&& f) &;
template<class F>
  constexpr auto and_then(F&& f) const&;
\end{itemdecl}

\begin{itemdescr}
\pnum
\constraints
\tcode{is_constructible_v<E, decltype(error())>} is \tcode{true}.

\pnum
\mandates
\tcode{remove_cvref_t<invoke_result_t<F>>} is a specialization of \tcode{expected} and its \tcode{error_type} is the same type as \tcode{E}.

\pnum
\effects
Equivalent to: \tcode{if (has_value()) return invoke(std::forward<F>(f)); else return U(unexpect, error());} where \tcode{U} is \tcode{remove_cvref_t<invoke_result_t<F>>}.
\end{itemdescr}

\indexlibrarymember{and_then}{expected}%
\begin{itemdecl}
template<class F>
  constexpr auto and_then(F&& f) &&;
template<class F>
  constexpr auto and_then(F&& f) const&&;
\end{itemdecl}

\begin{itemdescr}
\pnum
\constraints
\tcode{is_constructible_v<E, decltype(std::move(error()))>} is \tcode{true}.

\pnum
\mandates
\tcode{remove_cvref_t<invoke_result_t<F>>} is a specialization of \tcode{expected} and its \tcode{error_type} is the same type as \tcode{E}.

\pnum
\effects
Equivalent to: \tcode{if (has_value()) return invoke(std::forward<F>(f)); else return U(unexpect, std::move(error()));} where \tcode{U} is \tcode{remove_cvref_t<invoke_result_t<F>>}.
\end{itemdescr}

\indexlibrarymember{or_else}{expected}%
\begin{itemdecl}
template<class F>
  constexpr auto or_else(F&& f) &;
template<class F>
  constexpr auto or_else(F&& f) const&;
\end{itemdecl}

\begin{itemdescr}
\pnum
\mandates
\tcode{remove_cvref_t<invoke_result_t<F, decltype(error())>>} is a specialization of \tcode{expected} and its \tcode{value_type} is the same type as \tcode{T}.

\pnum
\effects
Equivalent to: \tcode{if (has_value()) return G(); else return invoke(std::forward<F>(f), error());} where \tcode{G} is \tcode{remove_cvref_t<invoke_result_t<F, decltype(error())>>}.
\end{itemdescr}

\indexlibrarymember{or_else}{expected}%
\begin{itemdecl}
template<class F>
  constexpr auto or_else(F&& f) &&;
template<class F>
  constexpr auto or_else(F&& f) const&&;
\end{itemdecl}

\begin{itemdescr}
\pnum
\mandates
\tcode{remove_cvref_t<invoke_result_t<F, decltype(std::move(error()))>>} is a specialization of \tcode{expected} and its \tcode{value_type} is the same type as \tcode{T}.

\pnum
\effects
Equivalent to: \tcode{if (has_value()) return G(); else return invoke(std::forward<F>(f), std::move(error()));} where \tcode{G} is \tcode{remove_cvref_t<invoke_result_t<F, decltype(std::move(error()))>>}.
\end{itemdescr}

\indexlibrarymember{transform}{expected}%
\begin{itemdecl}
template<class F>
  constexpr auto transform(F&& f) &;
template<class F>
  constexpr auto transform(F&& f) const&;
\end{itemdecl}

\begin{itemdescr}
\pnum
\constraints
\tcode{is_constructible_v<E, decltype(error())>} is \tcode{true}.

\pnum
\mandates
\tcode{U} is a valid value type for \tcode{expected}, where \tcode{U} is \tcode{remove_cv_t<invoke_result_t<F>>}.

\pnum
\effects
If \tcode{has_value()} is \tcode{false}, returns \tcode{expected<U, E>(unexpect, error())}. Otherwise, if \tcode{is_void_v<U>} is \tcode{false}, returns an \tcode{expected<U, E>} object whose \exposid{has_val} member is \tcode{true} and \tcode{val} member is direct-non-list-initialized with \tcode{invoke(std::forward<F>(f))}. Otherwise, evaluates \tcode{invoke(std::forward<F>(f))} and then returns \tcode{expected<U, E>()}.
\end{itemdescr}

\indexlibrarymember{transform}{expected}%
\begin{itemdecl}
template<class F>
  constexpr auto transform(F&& f) &&;
template<class F>
  constexpr auto transform(F&& f) const&&;
\end{itemdecl}

\begin{itemdescr}
\pnum
\constraints
\tcode{is_constructible_v<E, decltype(std::move(error()))>} is \tcode{true}.

\pnum
\mandates
\tcode{U} is a valid value type for \tcode{expected}, where \tcode{U} is \tcode{remove_cv_t<invoke_result_t<F>>}.

\pnum
\effects
If \tcode{has_value()} is \tcode{false}, returns \tcode{expected<U, E>(unexpect, std::move(error()))}. Otherwise, if \tcode{is_void_v<U>} is \tcode{false}, returns an \tcode{expected<U, E>} object whose \exposid{has_val} member is \tcode{true} and \tcode{val} member is direct-non-list-initialized with \tcode{invoke(std::forward<F>(f))}. Otherwise, evaluates \tcode{invoke(std::forward<F>(f))} and then returns \tcode{expected<U, E>()}.
\end{itemdescr}

\indexlibrarymember{transform_error}{expected}%
\begin{itemdecl}
template<class F>
  constexpr auto transform_error(F&& f) &;
template<class F>
  constexpr auto transform_error(F&& f) const&;
\end{itemdecl}

\begin{itemdescr}
\pnum
\mandates
\tcode{G} is a valid template argument for \tcode{unexpected} and the declaration \tcode{G g(invoke(std::forward<F>(f), error()));} is well-formed, where \tcode{G} is \tcode{remove_cv_t<invoke_result_t<F, decltype(error())>>}.

\pnum
\returns
If \tcode{has_value()} is \tcode{true}, \tcode{expected<T, G>()}; otherwise, an \tcode{expected<T, G>} object whose \exposid{has_val} member is \tcode{false} and \exposid{unex} member is direct-non-list-initialized with \tcode{invoke(std::forward<F>(f), error())}.
\end{itemdescr}

\indexlibrarymember{transform_error}{expected}%
\begin{itemdecl}
template<class F>
  constexpr auto transform_error(F&& f) &&;
template<class F>
  constexpr auto transform_error(F&& f) const&&;
\end{itemdecl}

\begin{itemdescr}
\pnum
\mandates
\tcode{G} is a valid template argument for \tcode{unexpected} and the declaration \tcode{G g(invoke(std::forward<F>(f), std::move(error())));} is well-formed, where \tcode{G} is \tcode{remove_cv_t<invoke_result_t<F, decltype(std::move(error()))>>}.

\pnum
\returns
If \tcode{has_value()} is \tcode{true}, \tcode{expected<T, G>()}; otherwise, an \tcode{expected<T, G>} object whose \exposid{has_val} member is \tcode{false} and \exposid{unex} member is direct-non-list-initialized with \tcode{invoke(std::forward<F>(f), std::move(error()))}.
\end{itemdescr}

\rSec3[expected.void.eq]{Equality operators}

\indexlibraryglobal{operator==}%
\begin{itemdecl}
template<class T2, class E2>
  friend constexpr bool operator==(const expected& x, const expected<T2, E2>& y);
\end{itemdecl}

\begin{itemdescr}
\pnum
\constraints
\tcode{is_void_v<T2>} is \tcode{true}.

\pnum
\mandates
\tcode{is_void_v<T2>} is \tcode{true}. The expression \tcode{x.error() == y.error()} is well-formed and its result is convertible to \tcode{bool}.

\pnum
\returns
If \tcode{x.has_value() != y.has_value()}, \tcode{false}; otherwise, if \tcode{x.has_value()} is \tcode{true}, \tcode{true}; otherwise \tcode{x.error() == y.error()}.
\end{itemdescr}

\indexlibraryglobal{operator==}%
\begin{itemdecl}
template<class E2>
  friend constexpr bool operator==(const expected& x, const unexpected<E2>& e);
\end{itemdecl}

\begin{itemdescr}
\pnum
\mandates
The expression \tcode{x.error() == e.error()} is well-formed and its result is convertible to \tcode{bool}.

\pnum
\returns
\tcode{!x.has_value() \&\& static_cast<bool>(x.error() == e.error())}.
\end{itemdescr}

\rSec2[expected.ref]{Partial specialization of expected for reference types}

\rSec3[expected.ref.general]{General}

\begin{codeblock}
namespace std {
  template<class T, class E>
  class @\libglobal{expected}@<T&, E> {
  private:
    using @\exposid{error-value-type}@ = remove_cv_t<remove_reference_t<E>>; // \expos

  public:
    using @\libmember{value_type}{expected}@ = T&;
    using @\libmember{error_type}{expected}@ = E;
    using @\libmember{unexpected_type}{expected}@ = unexpected<E>;

    template<class U> using rebind = expected<U, error_type>;

    expected() = delete;

    constexpr expected(const expected&) = default;

    constexpr expected(const expected& rhs) noexcept(
        is_nothrow_copy_constructible_v<E>);

    constexpr expected(expected&&) = default;

    constexpr expected(expected&& rhs) noexcept(is_nothrow_move_constructible_v<E>);

    template<class... Args> constexpr expected(in_place_t, Args&&...) = delete;

    template<class U = T>
    constexpr explicit(!is_convertible_v<U, T&>)
        expected(U&& u) noexcept(is_nothrow_constructible_v<T&, U>);

    template<class U> constexpr expected(U&&) = delete;

    template<class U, class G>
    constexpr explicit(!is_convertible_v<U&, T&> || !is_convertible_v<const G&, E>)
        expected(const expected<U&, G>& rhs);

    template<class U, class G>
    constexpr explicit(!is_convertible_v<U&, T&> || !is_convertible_v<G, E>)
        expected(expected<U&, G>&& rhs);

    template<class U, class G>
    constexpr explicit(!is_convertible_v<U&, T&> || !is_convertible_v<G, E>)
        expected(const expected<U&, G>& rhs);

    template<class U, class G>
    constexpr explicit(!is_convertible_v<U&, T&> || !is_convertible_v<G, E>)
        expected(expected<U&, G>&& rhs);

    template<class G>
    constexpr explicit(!is_convertible_v<const G&, E>) expected(const unexpected<G>& e);

    template<class G>
    constexpr explicit(!is_convertible_v<G, E>) expected(unexpected<G>&& e);

    template<class G>
    constexpr explicit(!is_convertible_v<G, E>)
        expected(const unexpected<G>& e) noexcept;

    template<class G>
    constexpr explicit(!is_convertible_v<G, E>) expected(unexpected<G>&& e) noexcept;

    template<class G> constexpr expected(const unexpected<G>&) = delete;

    template<class G>
      requires(is_reference_v<E> && !is_reference_v<G>)
    constexpr expected(unexpected<G>&&) = delete;

    template<class... Args> constexpr explicit expected(unexpect_t, Args&&... args);

    template<class... Args> constexpr expected(unexpect_t, Args&&...) = delete;

    template<class... Args>
      requires(is_reference_v<E> && !is_constructible_v<E, Args...> &&
               !@\exposid{unexpect-dangles-v}@<E, Args...>)
    constexpr expected(unexpect_t, Args&&...) = delete;

    template<class U, class... Args>
    constexpr explicit expected(unexpect_t, initializer_list<U> il, Args&&... args);

    template<class U, class... Args>
    constexpr expected(unexpect_t, initializer_list<U>, Args&&...) = delete;

    constexpr ~expected() = default;

    constexpr ~expected();

    constexpr expected& operator=(const expected&) = default;

    constexpr expected& operator=(const expected& rhs) noexcept(
        is_nothrow_copy_constructible_v<E> && is_nothrow_copy_assignable_v<E>);

    constexpr expected& operator=(expected&&) noexcept = default;

    constexpr expected& operator=(expected&& rhs) noexcept(
        is_nothrow_move_constructible_v<E> && is_nothrow_move_assignable_v<E>);

    template<class U = T> constexpr expected& operator=(U&& u);

    template<class G> constexpr expected& operator=(const unexpected<G>& e);

    template<class G> constexpr expected& operator=(unexpected<G>&& e);

    template<class G> constexpr expected& operator=(const unexpected<G>& e);

    template<class G> constexpr expected& operator=(unexpected<G>&& e);

    template<class G> constexpr expected& operator=(const unexpected<G>&) = delete;

    template<class G>
      requires(is_reference_v<E> && !is_reference_v<G>)
    constexpr expected& operator=(unexpected<G>&&) = delete;

    template<class U = T>
    constexpr T& emplace(U&& u) noexcept(is_nothrow_constructible_v<T&, U>);

    constexpr void swap(expected& rhs) noexcept(is_nothrow_move_constructible_v<E> &&
                                                (is_reference_v<E> ||
                                                 is_nothrow_swappable_v<E>));

    friend constexpr void swap(expected& x, expected& y) noexcept(noexcept(x.swap(y)));

    constexpr T* operator->() const noexcept;
    constexpr T& operator*() const noexcept;

    constexpr explicit operator bool() const noexcept;
    constexpr bool has_value() const noexcept;

    constexpr T& value() const&;
    constexpr T& value() &&;

    constexpr const E& error() const& noexcept;
    constexpr E& error() & noexcept;
    constexpr const E&& error() const&& noexcept;
    constexpr E&& error() && noexcept;

    template<class U = remove_cv_t<T>> constexpr remove_cv_t<T> value_or(U&& def) const;

    template<class G = @\exposid{error-value-type}@>
    constexpr @\exposid{error-value-type}@ error_or(G&& def) const&;

    template<class G = @\exposid{error-value-type}@>
    constexpr @\exposid{error-value-type}@ error_or(G&& def) &&;

    template<class F> constexpr auto and_then(F&& f) &;
    template<class F> constexpr auto and_then(F&& f) &&;
    template<class F> constexpr auto and_then(F&& f) const&;
    template<class F> constexpr auto and_then(F&& f) const&&;

    template<class F> constexpr auto or_else(F&& f) &;
    template<class F> constexpr auto or_else(F&& f) &&;
    template<class F> constexpr auto or_else(F&& f) const&;
    template<class F> constexpr auto or_else(F&& f) const&&;

    template<class F> constexpr auto transform(F&& f) &;
    template<class F> constexpr auto transform(F&& f) &&;
    template<class F> constexpr auto transform(F&& f) const&;
    template<class F> constexpr auto transform(F&& f) const&&;

    template<class F> constexpr auto transform_error(F&& f) &;
    template<class F> constexpr auto transform_error(F&& f) &&;
    template<class F> constexpr auto transform_error(F&& f) const&;
    template<class F> constexpr auto transform_error(F&& f) const&&;

    template<class T2, class E2>
    friend constexpr bool operator==(const expected& x, const expected<T2, E2>& y);

    template<class T2>
    friend constexpr bool operator==(const expected& x, const T2& val);

    template<class E2>
    friend constexpr bool operator==(const expected& x, const unexpected<E2>& e);

  private:
    bool @\exposid{has_val}@; // \expos
    union {
      T* @\exposid{val}@;             // \expos
      unexpected<E> @\exposid{unex}@; // \expos
    };
  };
}
\end{codeblock}

\pnum
\mandates
A program that instantiates the definition of \tcode{expected<T\&, E>} with an \tcode{E} that is not a valid template argument for \tcode{unexpected} is ill-formed. \tcode{T} shall be an object type that is not an array type, \tcode{in_place_t}, \tcode{unexpect_t}, or a specialization of \tcode{unexpected}.

\pnum
\remarks
An object of type \tcode{expected<T\&, E>} either represents a reference to an object of type \tcode{T}, or holds an error. Member \exposid{has_val} indicates whether the object represents a reference. When it represents a reference, member \exposid{val} points to the referenced object, which is not owned by the \tcode{expected} object. Otherwise, the error is \tcode{\exposid{unex}.error()}.

\rSec3[expected.ref.cons]{Constructors}

\indexlibraryctor{expected}%
\begin{itemdecl}
expected() = delete;
\end{itemdecl}

\begin{itemdescr}
\pnum
\remarks
\tcode{expected<T\&, E>} has no default constructor: a reference cannot be null, so there is no empty state to default-construct into.
\end{itemdescr}

\indexlibraryctor{expected}%
\begin{itemdecl}
constexpr expected(const expected&) = default;
\end{itemdecl}

\begin{itemdescr}
\pnum
\effects
If \tcode{rhs.has_value()} is \tcode{true}, initializes \exposid{val} with \tcode{rhs.\exposid{val}}, so that \tcode{*this} and \tcode{rhs} refer to the same object; otherwise, initializes \exposid{unex} with \tcode{rhs.\exposid{unex}}.

\pnum
\ensures
\tcode{rhs.has_value() == this->has_value()}.

\pnum
\remarks
This constructor is trivial.
\end{itemdescr}

\indexlibraryctor{expected}%
\begin{itemdecl}
constexpr expected(const expected& rhs) noexcept(is_nothrow_copy_constructible_v<E>);
constexpr expected(expected&& rhs) noexcept(is_nothrow_move_constructible_v<E>);
\end{itemdecl}

\begin{itemdescr}
\pnum
\constraints
\tcode{is_copy_constructible_v<E>} is \tcode{true} and \tcode{is_trivially_copy_constructible_v<E>} is \tcode{false}.

\pnum
\effects
If \tcode{rhs.has_value()} is \tcode{true}, initializes \exposid{val} with \tcode{rhs.\exposid{val}}, so that \tcode{*this} and \tcode{rhs} refer to the same object; otherwise, initializes \exposid{unex} with \tcode{rhs.\exposid{unex}}.

\pnum
\ensures
\tcode{rhs.has_value() == this->has_value()}.

\pnum
\remarks
The exception specification of the move constructor is equivalent to \tcode{is_nothrow_move_constructible_v<E>}. Each of these constructors is trivial if the corresponding constructor of \tcode{E} is trivial. The copy constructor is defined as deleted unless \tcode{is_copy_constructible_v<E>} is \tcode{true}, and the move constructor unless \tcode{is_move_constructible_v<E>} is \tcode{true}.
\end{itemdescr}

\indexlibraryctor{expected}%
\begin{itemdecl}
constexpr expected(expected&&) = default;
\end{itemdecl}

\begin{itemdescr}
\pnum
\effects
If \tcode{rhs.has_value()} is \tcode{true}, initializes \exposid{val} with \tcode{rhs.\exposid{val}}, so that \tcode{*this} and \tcode{rhs} refer to the same object; otherwise, initializes \exposid{unex} with \tcode{std::move(rhs.\exposid{unex})}.

\pnum
\ensures
\tcode{rhs.has_value() == this->has_value()}.

\pnum
\remarks
This constructor is trivial.
\end{itemdescr}

\indexlibraryctor{expected}%
\begin{itemdecl}
template<class... Args>
  constexpr expected(in_place_t, Args&&...) = delete;
\end{itemdecl}

\begin{itemdescr}
\pnum
\remarks
\tcode{expected<T\&, E>} has no in-place value constructor: \tcode{T\&} cannot be constructed in-place. Pass a \tcode{U} convertible to \tcode{T\&} instead.
\end{itemdescr}

\indexlibraryctor{expected}%
\begin{itemdecl}
template<class U = T>
  constexpr explicit(!is_convertible_v<U, T&>) expected(U&& u) noexcept(is_nothrow_constructible_v<T&, U>);
\end{itemdecl}

\begin{itemdescr}
\pnum
\constraints
\tcode{remove_cvref_t<U>} is not \tcode{in_place_t}, \tcode{expected}, or a specialization of \tcode{unexpected}; \tcode{is_constructible_v<T\&, U>} is \tcode{true}; and \tcode{reference_constructs_from_temporary_v<T\&, U>} is \tcode{false}.

\pnum
\effects
Let \tcode{r} be the lvalue result of \tcode{T\& r = std::forward<U>(u);}. Initializes \exposid{val} with \tcode{addressof(r)}.

\pnum
\ensures
\tcode{has_value()} is \tcode{true}.
\end{itemdescr}

\indexlibraryctor{expected}%
\begin{itemdecl}
template<class U>
  constexpr expected(U&&) = delete;
\end{itemdecl}

\begin{itemdescr}
\pnum
\constraints
\tcode{reference_constructs_from_temporary_v<T\&, U>} is \tcode{true}.

\pnum
\remarks
A constructor for which \tcode{reference_constructs_from_temporary_v<T\&, U>} is \tcode{true} — one that would bind \tcode{T\&} to a temporary — is defined as deleted.
\end{itemdescr}

\indexlibraryctor{expected}%
\begin{itemdecl}
template<class U, class G>
  constexpr explicit(!is_convertible_v<U&, T&> || !is_convertible_v<const G&, E>) expected(const expected<U&, G>& rhs);
template<class U, class G>
  constexpr explicit(!is_convertible_v<U&, T&> || !is_convertible_v<G, E>) expected(expected<U&, G>&& rhs);
\end{itemdecl}

\begin{itemdescr}
\pnum
\constraints
\tcode{is_constructible_v<T\&, U\&>} is \tcode{true}; \tcode{reference_constructs_from_temporary_v<T\&, U\&>} is \tcode{false}; and \tcode{is_constructible_v<E, const G\&>} is \tcode{true}.

\pnum
\effects
If \tcode{rhs.has_value()} is \tcode{true}, initializes \exposid{val} with \tcode{addressof(*rhs)}, so that \tcode{*this} refers to the object referred to by \tcode{rhs}; otherwise, initializes \exposid{unex} with the error of \tcode{rhs}. No object referred to by \tcode{rhs} is moved from.
\end{itemdescr}

\indexlibraryctor{expected}%
\begin{itemdecl}
template<class U, class G>
  constexpr explicit(!is_convertible_v<U&, T&> || !is_convertible_v<G, E>) expected(const expected<U&, G>& rhs);
template<class U, class G>
  constexpr explicit(!is_convertible_v<U&, T&> || !is_convertible_v<G, E>) expected(expected<U&, G>&& rhs);
\end{itemdecl}

\begin{itemdescr}
\pnum
\constraints
\tcode{is_constructible_v<T\&, U\&>} is \tcode{true}; \tcode{reference_constructs_from_temporary_v<T\&, U\&>} is \tcode{false}; \tcode{G} is a reference type; and \tcode{is_convertible_v<G, E>} is \tcode{true}.

\pnum
\effects
If \tcode{rhs.has_value()} is \tcode{true}, initializes \exposid{val} with \tcode{addressof(*rhs)}, so that \tcode{*this} refers to the object referred to by \tcode{rhs}; otherwise, initializes \exposid{unex} with the error of \tcode{rhs}. No object referred to by \tcode{rhs} is moved from.

\pnum
\remarks
Unlike the value-\tcode{E} overload above, this overload participates in overload resolution only when \tcode{E} and \tcode{G} are both reference types, so the referenced error object is never copied.
\end{itemdescr}

\indexlibraryctor{expected}%
\begin{itemdecl}
template<class G>
  constexpr explicit(!is_convertible_v<const G&, E>) expected(const unexpected<G>& e);
template<class G>
  constexpr explicit(!is_convertible_v<G, E>) expected(unexpected<G>&& e);
\end{itemdecl}

\begin{itemdescr}
\pnum
\constraints
\tcode{is_constructible_v<E, const G\&>} is \tcode{true}.

\pnum
\effects
Initializes \exposid{unex} with the error of \tcode{e}.

\pnum
\ensures
\tcode{has_value()} is \tcode{false}.
\end{itemdescr}

\indexlibraryctor{expected}%
\begin{itemdecl}
template<class G>
  constexpr explicit(!is_convertible_v<G, E>) expected(const unexpected<G>& e) noexcept;
template<class G>
  constexpr explicit(!is_convertible_v<G, E>) expected(unexpected<G>&& e) noexcept;
\end{itemdecl}

\begin{itemdescr}
\pnum
\constraints
\tcode{is_reference_v<G>} is \tcode{true}; \tcode{is_constructible_v<E, G>} is \tcode{true}; and \tcode{reference_constructs_from_temporary_v<E, G>} is \tcode{false}.

\pnum
\effects
Initializes \exposid{unex} with the error of \tcode{e}.

\pnum
\ensures
\tcode{has_value()} is \tcode{false}.

\pnum
\remarks
This constructor never throws: the referent is bound, not copied.
\end{itemdescr}

\indexlibraryctor{expected}%
\begin{itemdecl}
template<class G>
  constexpr expected(const unexpected<G>&) = delete;
\end{itemdecl}

\begin{itemdescr}
\pnum
\constraints
\tcode{is_reference_v<E>} is \tcode{true} and \tcode{is_reference_v<G>} is \tcode{false}.

\pnum
\remarks
When \tcode{E} is a reference type, an overload taking \tcode{unexpected<G>} for a non-reference \tcode{G} is defined as deleted: the referent would live inside the (possibly temporary) source \tcode{unexpected<G>} object, and binding \tcode{E\&} to it would dangle. Use an \tcode{unexpected<E\&>} holding an external object instead.
\end{itemdescr}

\indexlibraryctor{expected}%
\begin{itemdecl}
template<class... Args>
  constexpr explicit expected(unexpect_t, Args&&... args);
\end{itemdecl}

\begin{itemdescr}
\pnum
\constraints
\tcode{is_constructible_v<E, Args...>} is \tcode{true}.

\pnum
\effects
Direct-non-list-initializes \exposid{unex} with \tcode{in_place} and \tcode{std::forward<Args>(args)...}.

\pnum
\ensures
\tcode{has_value()} is \tcode{false}.
\end{itemdescr}

\indexlibraryctor{expected}%
\begin{itemdecl}
template<class... Args>
  constexpr expected(unexpect_t, Args&&...) = delete;
\end{itemdecl}

\begin{itemdescr}
\pnum
\constraints
\tcode{\exposid{unexpect-dangles-v}<E, Args...>} is \tcode{true}.

\pnum
\remarks
When \tcode{E} is a reference type, an overload with the same parameter types is defined as deleted if the single argument would bind \tcode{E\&} to a temporary, or if it is otherwise not usable to construct \tcode{E}.
\end{itemdescr}

\indexlibraryctor{expected}%
\begin{itemdecl}
template<class U, class... Args>
  constexpr explicit expected(unexpect_t, initializer_list<U> il, Args&&... args);
\end{itemdecl}

\begin{itemdescr}
\pnum
\constraints
\tcode{is_reference_v<E>} is \tcode{false}, and \tcode{is_constructible_v<E, initializer_list<U>\&, Args...>} is \tcode{true}.

\pnum
\effects
Direct-non-list-initializes \exposid{unex} with \tcode{in_place}, \tcode{il}, and \tcode{std::forward<Args>(args)...}.

\pnum
\ensures
\tcode{has_value()} is \tcode{false}.

\pnum
\remarks
An overload with the same parameter types is defined as deleted when \tcode{E} is an lvalue reference type. An initializer list cannot provide the required long-lived error referent.
\end{itemdescr}

\indexlibraryctor{expected}%
\begin{itemdecl}
template<class U, class... Args>
  constexpr expected(unexpect_t, initializer_list<U>, Args&&...) = delete;
\end{itemdecl}

\begin{itemdescr}
\pnum
\constraints
\tcode{is_reference_v<E>} is \tcode{true}.

\pnum
\remarks
An overload with the same parameter types is defined as deleted when \tcode{E} is an lvalue reference type. An initializer list cannot provide the required long-lived error referent.
\end{itemdescr}

\rSec3[expected.ref.dtor]{Destructor}

\indexlibrarydtor{expected}%
\begin{itemdecl}
constexpr ~expected();
\end{itemdecl}

\begin{itemdescr}
\pnum
\constraints
\tcode{is_trivially_destructible_v<E>} is \tcode{true}.

\pnum
\effects
None: \tcode{*this} never owns the object it refers to; \tcode{T} is never destroyed.

\pnum
\remarks
This destructor is trivial.
\end{itemdescr}

\indexlibrarydtor{expected}%
\begin{itemdecl}
constexpr ~expected();
\end{itemdecl}

\begin{itemdescr}
\pnum
\constraints
\tcode{is_trivially_destructible_v<E>} is \tcode{false}.

\pnum
\effects
If \tcode{has_value()} is \tcode{false}, destroys \exposid{unex}. \tcode{T} is not destroyed; \tcode{*this} never owns the object it refers to.

\pnum
\remarks
This destructor is trivial if \tcode{E} is trivially destructible.
\end{itemdescr}

\rSec3[expected.ref.assign]{Assignment}

\indexlibrarymember{operator=}{expected}%
\begin{itemdecl}
constexpr expected& operator=(const expected&);
\end{itemdecl}

\begin{itemdescr}
\pnum
\constraints
\tcode{is_trivially_copy_constructible_v<E>} is \tcode{true}, \tcode{is_trivially_copy_assignable_v<E>} is \tcode{true}, and \tcode{is_trivially_destructible_v<E>} is \tcode{true}.

\pnum
\effects
If \tcode{rhs.has_value()} is \tcode{true}: if \tcode{has_value()} is \tcode{true}, assigns \tcode{rhs.\exposid{val}} to \exposid{val}; otherwise destroys \exposid{unex} and initializes \exposid{val} with \tcode{rhs.\exposid{val}}. If \tcode{rhs.has_value()} is \tcode{false}, the error of \tcode{rhs} is assigned to or used to initialize \exposid{unex}, as for the primary template. In every case \tcode{*this} comes to refer to the object \tcode{rhs} refers to, or to hold the error of \tcode{rhs}.

\pnum
\returns
\tcode{*this}.

\pnum
\remarks
Assignment rebinds: assigning to an \tcode{expected<T\&, E>} that holds a value changes which object it refers to. It never assigns through to the referent. This operator is trivial.
\end{itemdescr}

\indexlibrarymember{operator=}{expected}%
\begin{itemdecl}
constexpr expected& operator=(const expected& rhs) noexcept(is_nothrow_copy_constructible_v<E> &&
                                                            is_nothrow_copy_assignable_v<E>);
constexpr expected& operator=(expected&& rhs) noexcept(is_nothrow_move_constructible_v<E> &&
                                                       is_nothrow_move_assignable_v<E>);
\end{itemdecl}

\begin{itemdescr}
\pnum
\constraints
\tcode{is_reference_v<E> || (is_copy_constructible_v<E> \&\& is_copy_assignable_v<E>)} is \tcode{true} and \tcode{(is_trivially_copy_constructible_v<E> \&\& is_trivially_copy_assignable_v<E> \&\& is_trivially_destructible_v<E>)} is \tcode{false}.

\pnum
\effects
If \tcode{rhs.has_value()} is \tcode{true}: if \tcode{has_value()} is \tcode{true}, assigns \tcode{rhs.\exposid{val}} to \exposid{val}; otherwise destroys \exposid{unex} and initializes \exposid{val} with \tcode{rhs.\exposid{val}}. If \tcode{rhs.has_value()} is \tcode{false}, the error of \tcode{rhs} is assigned to or used to initialize \exposid{unex}, as for the primary template. In every case \tcode{*this} comes to refer to the object \tcode{rhs} refers to, or to hold the error of \tcode{rhs}.

\pnum
\returns
\tcode{*this}.

\pnum
\remarks
The exception specification of the move assignment operator is equivalent to \tcode{is_nothrow_move_constructible_v<E> \&\& is_nothrow_move_assignable_v<E>}. The copy assignment operator is defined as deleted unless \tcode{is_copy_assignable_v<E>} is \tcode{true} and \tcode{is_copy_constructible_v<E>} is \tcode{true}, and the move assignment operator unless \tcode{is_move_assignable_v<E>} is \tcode{true} and \tcode{is_move_constructible_v<E>} is \tcode{true}.
\end{itemdescr}

\indexlibrarymember{operator=}{expected}%
\begin{itemdecl}
constexpr expected& operator=(expected&&) noexcept;
\end{itemdecl}

\begin{itemdescr}
\pnum
\constraints
\tcode{is_trivially_move_constructible_v<E>} is \tcode{true}, \tcode{is_trivially_move_assignable_v<E>} is \tcode{true}, and \tcode{is_trivially_destructible_v<E>} is \tcode{true}.

\pnum
\effects
If \tcode{rhs.has_value()} is \tcode{true}: if \tcode{has_value()} is \tcode{true}, assigns \tcode{rhs.\exposid{val}} to \exposid{val}; otherwise destroys \exposid{unex} and initializes \exposid{val} with \tcode{rhs.\exposid{val}}. If \tcode{rhs.has_value()} is \tcode{false}, the error of \tcode{rhs} is assigned to or used to initialize \exposid{unex}, as for the primary template.

\pnum
\returns
\tcode{*this}.

\pnum
\remarks
This operator is trivial.
\end{itemdescr}

\indexlibrarymember{operator=}{expected}%
\begin{itemdecl}
template<class U = T>
  constexpr expected& operator=(U&& u);
\end{itemdecl}

\begin{itemdescr}
\pnum
\constraints
\tcode{remove_cvref_t<U>} is neither \tcode{expected} nor a specialization of \tcode{unexpected}, \tcode{is_constructible_v<T\&, U>} is \tcode{true}, and \tcode{reference_constructs_from_temporary_v<T\&, U>} is \tcode{false}.

\pnum
\effects
Let \tcode{r} be the lvalue result of \tcode{T\& r = std::forward<U>(u);}. If \tcode{has_value()} is \tcode{true}, assigns \tcode{addressof(r)} to \exposid{val}. Otherwise, destroys \exposid{unex}, initializes \exposid{val} with \tcode{addressof(r)}, and sets \exposid{has_val} to \tcode{true}; if binding \tcode{r} throws, \tcode{*this} is left unchanged.

\pnum
\returns
\tcode{*this}.
\end{itemdescr}

\indexlibrarymember{operator=}{expected}%
\begin{itemdecl}
template<class G>
  constexpr expected& operator=(const unexpected<G>& e);
template<class G>
  constexpr expected& operator=(unexpected<G>&& e);
\end{itemdecl}

\begin{itemdescr}
\pnum
\constraints
\tcode{is_constructible_v<E, const G\&>} is \tcode{true} and \tcode{is_assignable_v<E\&, const G\&>} is \tcode{true}.

\pnum
\effects
Makes \tcode{*this} hold the error of \tcode{e}, reinitializing \exposid{unex} from \tcode{e} rather than assigning through it.

\pnum
\returns
\tcode{*this}.
\end{itemdescr}

\indexlibrarymember{operator=}{expected}%
\begin{itemdecl}
template<class G>
  constexpr expected& operator=(const unexpected<G>& e);
template<class G>
  constexpr expected& operator=(unexpected<G>&& e);
\end{itemdecl}

\begin{itemdescr}
\pnum
\constraints
\tcode{is_reference_v<G>} is \tcode{true}; \tcode{is_constructible_v<E, G>} is \tcode{true}; and \tcode{reference_constructs_from_temporary_v<E, G>} is \tcode{false}.

\pnum
\effects
Makes \tcode{*this} hold the error of \tcode{e}, reinitializing \exposid{unex} from \tcode{e} rather than assigning through it. \tcode{\exposid{unex}.error()} thereafter refers to the same object as \tcode{e.error()}; the previously referenced object, if any, is not modified.

\pnum
\returns
\tcode{*this}.

\pnum
\remarks
This operator never throws: the referent is bound, not copied.
\end{itemdescr}

\indexlibrarymember{operator=}{expected}%
\begin{itemdecl}
template<class G>
  constexpr expected& operator=(const unexpected<G>&) = delete;
\end{itemdecl}

\begin{itemdescr}
\pnum
\constraints
\tcode{is_reference_v<E>} is \tcode{true} and \tcode{is_reference_v<G>} is \tcode{false}.

\pnum
\remarks
When \tcode{E} is a reference type, an overload taking \tcode{unexpected<G>} for a non-reference \tcode{G} is defined as deleted: it would rebind \exposid{unex} to \tcode{unexpected<G>}'s temporary storage.
\end{itemdescr}

\indexlibrarymember{emplace}{expected}%
\begin{itemdecl}
template<class U = T>
  constexpr T& emplace(U&& u) noexcept(is_nothrow_constructible_v<T&, U>);
\end{itemdecl}

\begin{itemdescr}
\pnum
\constraints
\tcode{is_constructible_v<T\&, U>} is \tcode{true} and \tcode{reference_constructs_from_temporary_v<T\&, U>} is \tcode{false}.

\pnum
\effects
Rebinds \tcode{*this} to refer to the object bound by \tcode{T\& r = std::forward<U>(u);}: if \tcode{has_value()} is \tcode{false}, destroys \exposid{unex} first. Sets \exposid{val} to \tcode{addressof(r)} and \exposid{has_val} to \tcode{true}.

\pnum
\returns
\tcode{*\exposid{val}}.
\end{itemdescr}

\rSec3[expected.ref.swap]{Swap}

\indexlibrarymember{swap}{expected}%
\begin{itemdecl}
constexpr void swap(expected& rhs) noexcept(is_nothrow_move_constructible_v<E> &&
                                            (is_reference_v<E> || is_nothrow_swappable_v<E>));
\end{itemdecl}

\begin{itemdescr}
\pnum
\constraints
\tcode{is_reference_v<E> || is_swappable_v<E>} is \tcode{true} and \tcode{is_move_constructible_v<E>} is \tcode{true}.

\pnum
\effects
Exchanges the states of \tcode{*this} and \tcode{rhs}. When both hold values, exchanges \exposid{val} and \tcode{rhs.\exposid{val}} — the referenced objects are not swapped. Otherwise behaves as the primary template's \tcode{swap} does for the error.

\pnum
\remarks
The exception specification is equivalent to \tcode{is_nothrow_move_constructible_v<E> \&\& (is_reference_v<E> || is_nothrow_swappable_v<E>)}.
\end{itemdescr}

\indexlibraryglobal{swap}%
\begin{itemdecl}
friend constexpr void swap(expected& x, expected& y) noexcept(noexcept(x.swap(y)));
\end{itemdecl}

\begin{itemdescr}
\pnum
\constraints
\tcode{is_reference_v<E> || is_swappable_v<E>} is \tcode{true} and \tcode{is_move_constructible_v<E>} is \tcode{true}.

\pnum
\effects
Equivalent to \tcode{x.swap(y)}.
\end{itemdescr}

\rSec3[expected.ref.obs]{Observers}

\indexlibrarymember{operator->}{expected}%
\begin{itemdecl}
constexpr T* operator->() const noexcept;
\end{itemdecl}

\begin{itemdescr}
\pnum
\expects
\tcode{has_value()} is \tcode{true}.

\pnum
\returns
\exposid{val}.

\pnum
\remarks
This is a \tcode{const} member function that returns a non-\tcode{const} \tcode{T*}; the constness of \tcode{*this} does not propagate to the referenced object. For deep \tcode{const}, use \tcode{expected<const T\&, E>}.
\end{itemdescr}

\indexlibrarymember{operator*}{expected}%
\begin{itemdecl}
constexpr T& operator*() const noexcept;
\end{itemdecl}

\begin{itemdescr}
\pnum
\expects
\tcode{has_value()} is \tcode{true}.

\pnum
\returns
\tcode{*\exposid{val}}.

\pnum
\remarks
This is a \tcode{const} member function that returns a non-\tcode{const} \tcode{T\&}; the constness of \tcode{*this} does not propagate to the referenced object. For deep \tcode{const}, use \tcode{expected<const T\&, E>}.
\end{itemdescr}

\indexlibrarymember{operator bool}{expected}%
\indexlibrarymember{has_value}{expected}%
\begin{itemdecl}
constexpr explicit operator bool() const noexcept;
constexpr bool has_value() const noexcept;
\end{itemdecl}

\begin{itemdescr}
\pnum
\returns
\exposid{has_val}.
\end{itemdescr}

\indexlibrarymember{value}{expected}%
\begin{itemdecl}
constexpr T& value() const&;
\end{itemdecl}

\begin{itemdescr}
\pnum
\mandates
\tcode{is_copy_constructible_v<\exposid{error-value-type}>} is \tcode{true}.

\pnum
\returns
\tcode{*\exposid{val}} if \tcode{has_value()} is \tcode{true}.

\pnum
\throws
\tcode{bad_expected_access(as_const(error()))} if \tcode{has_value()} is \tcode{false}.
\end{itemdescr}

\indexlibrarymember{value}{expected}%
\begin{itemdecl}
constexpr T& value() &&;
\end{itemdecl}

\begin{itemdescr}
\pnum
\returns
\tcode{*\exposid{val}} if \tcode{has_value()} is \tcode{true}.

\pnum
\throws
\tcode{bad_expected_access(std::move(error()))} if \tcode{has_value()} is \tcode{false}.
\end{itemdescr}

\indexlibrarymember{error}{expected}%
\begin{itemdecl}
constexpr const E& error() const& noexcept;
constexpr E& error() & noexcept;
\end{itemdecl}

\begin{itemdescr}
\pnum
\expects
\tcode{has_value()} is \tcode{false}.

\pnum
\returns
\tcode{\exposid{unex}.error()}.
\end{itemdescr}

\indexlibrarymember{error}{expected}%
\begin{itemdecl}
constexpr const E&& error() const&& noexcept;
constexpr E&& error() && noexcept;
\end{itemdecl}

\begin{itemdescr}
\pnum
\expects
\tcode{has_value()} is \tcode{false}.

\pnum
\returns
\tcode{std::move(\exposid{unex}).error()}.
\end{itemdescr}

\indexlibrarymember{value_or}{expected}%
\begin{itemdecl}
template<class U = remove_cv_t<T>>
  constexpr remove_cv_t<T> value_or(U&& def) const;
\end{itemdecl}

\begin{itemdescr}
\pnum
\constraints
\tcode{is_object_v<T>} is \tcode{true} and \tcode{is_array_v<T>} is \tcode{false}.

\pnum
\mandates
\tcode{is_convertible_v<T\&, remove_cv_t<T>>} and \tcode{is_convertible_v<U, remove_cv_t<T>>} are \tcode{true}.

\pnum
\returns
\tcode{has_value() ? static_cast<remove_cv_t<T>>(*\exposid{val}) : static_cast<remove_cv_t<T>>(std::forward<U>(def))}. The result is an object, never a reference.
\end{itemdescr}

\indexlibrarymember{error_or}{expected}%
\begin{itemdecl}
template<class G = @\exposid{error-value-type}@>
  constexpr @\exposid{error-value-type}@ error_or(G&& def) const&;
\end{itemdecl}

\begin{itemdescr}
\pnum
\constraints
\tcode{is_copy_constructible_v<remove_cv_t<remove_reference_t<E>>>} is \tcode{true} and \tcode{is_convertible_v<G, remove_cv_t<remove_reference_t<E>>>} is \tcode{true}.

\pnum
\returns
\tcode{std::forward<G>(def)} if \tcode{has_value()} is \tcode{true}, \tcode{error()} otherwise. The result is an object, never a reference.
\end{itemdescr}

\indexlibrarymember{error_or}{expected}%
\begin{itemdecl}
template<class G = @\exposid{error-value-type}@>
  constexpr @\exposid{error-value-type}@ error_or(G&& def) &&;
\end{itemdecl}

\begin{itemdescr}
\pnum
\constraints
\tcode{is_move_constructible_v<remove_cv_t<remove_reference_t<E>>>} is \tcode{true} and \tcode{is_convertible_v<G, remove_cv_t<remove_reference_t<E>>>} is \tcode{true}.

\pnum
\returns
\tcode{std::forward<G>(def)} if \tcode{has_value()} is \tcode{true}, \tcode{std::move(error())} otherwise. The result is an object, never a reference.
\end{itemdescr}

\rSec3[expected.ref.monadic]{Monadic operations}

\indexlibrarymember{and_then}{expected}%
\begin{itemdecl}
template<class F>
  constexpr auto and_then(F&& f) &;
template<class F>
  constexpr auto and_then(F&& f) const&;
\end{itemdecl}

\begin{itemdescr}
\pnum
\constraints
\tcode{is_constructible_v<E, decltype(error())>} is \tcode{true}.

\pnum
\mandates
\tcode{remove_cvref_t<invoke_result_t<F, T\&>>} is a specialization of \tcode{expected} and its \tcode{error_type} is the same type as \tcode{E}.

\pnum
\effects
Equivalent to: \tcode{if (has_value()) return invoke(std::forward<F>(f), *\exposid{val}); else return U(unexpect, error());} where \tcode{U} is \tcode{remove_cvref_t<invoke_result_t<F, T\&>>}.

\pnum
\remarks
The member templates \tcode{and_then}, \tcode{or_else}, \tcode{transform}, and \tcode{transform_error} behave as specified for the primary template, with one difference: the value is passed to the callable as \tcode{T\&} for every ref-qualification of \tcode{*this}. An rvalue \tcode{expected<T\&, E>} does not pass its referent as an rvalue; the object referred to is never moved from by these operations.
\end{itemdescr}

\indexlibrarymember{and_then}{expected}%
\begin{itemdecl}
template<class F>
  constexpr auto and_then(F&& f) &&;
template<class F>
  constexpr auto and_then(F&& f) const&&;
\end{itemdecl}

\begin{itemdescr}
\pnum
\constraints
\tcode{is_constructible_v<E, decltype(std::move(error()))>} is \tcode{true}.

\pnum
\mandates
\tcode{remove_cvref_t<invoke_result_t<F, T\&>>} is a specialization of \tcode{expected} and its \tcode{error_type} is the same type as \tcode{E}.

\pnum
\effects
Equivalent to: \tcode{if (has_value()) return invoke(std::forward<F>(f), *\exposid{val}); else return U(unexpect, std::move(error()));} where \tcode{U} is \tcode{remove_cvref_t<invoke_result_t<F, T\&>>}.
\end{itemdescr}

\indexlibrarymember{or_else}{expected}%
\begin{itemdecl}
template<class F>
  constexpr auto or_else(F&& f) &;
template<class F>
  constexpr auto or_else(F&& f) const&;
\end{itemdecl}

\begin{itemdescr}
\pnum
\mandates
\tcode{remove_cvref_t<invoke_result_t<F, decltype(error())>>} is a specialization of \tcode{expected} and its \tcode{value_type} is the same type as \tcode{T\&}.

\pnum
\effects
Equivalent to: \tcode{if (has_value()) return G(*\exposid{val}); else return invoke(std::forward<F>(f), error());} where \tcode{G} is \tcode{remove_cvref_t<invoke_result_t<F, decltype(error())>>}.
\end{itemdescr}

\indexlibrarymember{or_else}{expected}%
\begin{itemdecl}
template<class F>
  constexpr auto or_else(F&& f) &&;
template<class F>
  constexpr auto or_else(F&& f) const&&;
\end{itemdecl}

\begin{itemdescr}
\pnum
\mandates
\tcode{remove_cvref_t<invoke_result_t<F, decltype(std::move(error()))>>} is a specialization of \tcode{expected} and its \tcode{value_type} is the same type as \tcode{T\&}.

\pnum
\effects
Equivalent to: \tcode{if (has_value()) return G(*\exposid{val}); else return invoke(std::forward<F>(f), std::move(error()));} where \tcode{G} is \tcode{remove_cvref_t<invoke_result_t<F, decltype(std::move(error()))>>}.
\end{itemdescr}

\indexlibrarymember{transform}{expected}%
\begin{itemdecl}
template<class F>
  constexpr auto transform(F&& f) &;
template<class F>
  constexpr auto transform(F&& f) const&;
\end{itemdecl}

\begin{itemdescr}
\pnum
\constraints
\tcode{is_constructible_v<E, decltype(error())>} is \tcode{true}.

\pnum
\effects
Equivalent to: \tcode{if (!has_value()) return U(unexpect, error()); else return expected<U2, E>(in_place, invoke(std::forward<F>(f), *\exposid{val}));} where \tcode{U2} is \tcode{remove_cv_t<invoke_result_t<F, T\&>>} and \tcode{U} is \tcode{expected<U2, E>}.
\end{itemdescr}

\indexlibrarymember{transform}{expected}%
\begin{itemdecl}
template<class F>
  constexpr auto transform(F&& f) &&;
template<class F>
  constexpr auto transform(F&& f) const&&;
\end{itemdecl}

\begin{itemdescr}
\pnum
\constraints
\tcode{is_constructible_v<E, decltype(std::move(error()))>} is \tcode{true}.

\pnum
\effects
Equivalent to: \tcode{if (!has_value()) return U(unexpect, std::move(error())); else return expected<U2, E>(in_place, invoke(std::forward<F>(f), *\exposid{val}));} where \tcode{U2} is \tcode{remove_cv_t<invoke_result_t<F, T\&>>} and \tcode{U} is \tcode{expected<U2, E>}.
\end{itemdescr}

\indexlibrarymember{transform_error}{expected}%
\begin{itemdecl}
template<class F>
  constexpr auto transform_error(F&& f) &;
template<class F>
  constexpr auto transform_error(F&& f) const&;
\end{itemdecl}

\begin{itemdescr}
\pnum
\mandates
\begin{itemize}
\item \tcode{is_object_v<G>} is \tcode{true},
\item \tcode{is_array_v<G>} is \tcode{false},
\item \tcode{is_same_v<G, remove_cv_t<G>>} is \tcode{true},
\item \tcode{\exposid{is-unexpected-specialization}<G>::value} is \tcode{false}.
\end{itemize}

\pnum
\effects
Equivalent to: \tcode{if (has_value()) return G(*\exposid{val}); else return expected<T\&, G2>(unexpect, invoke(std::forward<F>(f), error()));} where \tcode{G2} is \tcode{remove_cv_t<invoke_result_t<F, decltype(error())>>} and \tcode{G} is \tcode{expected<T\&, G2>}.
\end{itemdescr}

\indexlibrarymember{transform_error}{expected}%
\begin{itemdecl}
template<class F>
  constexpr auto transform_error(F&& f) &&;
template<class F>
  constexpr auto transform_error(F&& f) const&&;
\end{itemdecl}

\begin{itemdescr}
\pnum
\mandates
\begin{itemize}
\item \tcode{is_object_v<G>} is \tcode{true},
\item \tcode{is_array_v<G>} is \tcode{false},
\item \tcode{is_same_v<G, remove_cv_t<G>>} is \tcode{true},
\item \tcode{\exposid{is-unexpected-specialization}<G>::value} is \tcode{false}.
\end{itemize}

\pnum
\effects
Equivalent to: \tcode{if (has_value()) return G(*\exposid{val}); else return expected<T\&, G2>(unexpect, invoke(std::forward<F>(f), std::move(error())));} where \tcode{G2} is \tcode{remove_cv_t<invoke_result_t<F, decltype(std::move(error()))>>} and \tcode{G} is \tcode{expected<T\&, G2>}.
\end{itemdescr}

\rSec3[expected.ref.eq]{Equality operators}

\indexlibraryglobal{operator==}%
\begin{itemdecl}
template<class T2, class E2>
  friend constexpr bool operator==(const expected& x, const expected<T2, E2>& y);
\end{itemdecl}

\begin{itemdescr}
\pnum
\constraints
\tcode{is_void_v<T2>} is \tcode{false}.

\pnum
\mandates
\tcode{!is_void_v<T2>} is \tcode{true}. The expression \tcode{*x == *y} is well-formed and its result is convertible to \tcode{bool}. The expression \tcode{x.error() == y.error()} is well-formed and its result is convertible to \tcode{bool}.

\pnum
\returns
If \tcode{x.has_value() != y.has_value()}, \tcode{false}; otherwise, if \tcode{x.has_value()} is \tcode{true}, \tcode{*x == *y}; otherwise \tcode{x.error() == y.error()}.

\pnum
\remarks
The equality operators behave as specified for the primary template, comparing referents through \tcode{operator*} and errors through \tcode{error()}.
\end{itemdescr}

\indexlibraryglobal{operator==}%
\begin{itemdecl}
template<class T2>
  friend constexpr bool operator==(const expected& x, const T2& val);
\end{itemdecl}

\begin{itemdescr}
\pnum
\constraints
\tcode{\exposid{is-expected-specialization}<T2>::value} is \tcode{false}.

\pnum
\mandates
\tcode{T2} is not a specialization of \tcode{expected}. The expression \tcode{*x == val} is well-formed and its result is convertible to \tcode{bool}.

\pnum
\returns
\tcode{x.has_value() \&\& static_cast<bool>(*x == val)}.
\end{itemdescr}

\indexlibraryglobal{operator==}%
\begin{itemdecl}
template<class E2>
  friend constexpr bool operator==(const expected& x, const unexpected<E2>& e);
\end{itemdecl}

\begin{itemdescr}
\pnum
\mandates
The expression \tcode{x.error() == e.error()} is well-formed and its result is convertible to \tcode{bool}.

\pnum
\returns
\tcode{!x.has_value() \&\& static_cast<bool>(x.error() == e.error())}.
\end{itemdescr}
